\section{Imported foundations}

% archon:covers MR4213770UniversalSecantBundlesAndSyzygiesOfCanonicalCurves/Foundations.lean

This chapter records the foundational cohomology-and-bundle layer on which every
later construction rests.  The substrate is supplied by \emph{Mathlib}: for a
scheme \(X\), the category \(X.\mathrm{Modules}\) of sheaves of
\(\cO_X\)-modules is abelian and carries an additive pushforward functor along a
morphism of schemes.  Producing the derived functors that compute sheaf
cohomology and higher direct images requires, in addition, that the relevant
category have injective resolutions, and here the two cohomological invariants
part ways into a \emph{light} and a \emph{heavy} route.

For \emph{scalar sheaf cohomology} \(H^i(X,\mathcal F)\) one only needs derived
functors in the category \(\mathrm{Sh}(X,\mathrm{Ab})\) of abelian-group sheaves
on the (essentially small) site of opens of \(X\).  That category \emph{is} of
the form \(\mathrm{Sheaf}\,J\,\mathcal A\) with \(\mathcal A=\mathrm{Ab}\), so
Mathlib's essentially-small theorem applies directly and makes it Grothendieck
abelian, hence equipped with injective resolutions, \emph{unconditionally}.  We
compute \(H^i\) by pushing an \(\cO_X\)-module to its underlying abelian-group
sheaf and deriving global sections there.

For \emph{higher direct images} \(R^if_*\) in the module category one genuinely
needs \(X.\mathrm{Modules}\) itself to be Grothendieck abelian.  This is the
\emph{heavy} bridge: sheaves of modules over a \emph{varying} ring sheaf are not
literally objects of a category \(\mathrm{Sheaf}\,J\,\mathcal A\), so the
essentially-small theorem does \emph{not} apply, and Grothendieck-abelianness
must be assembled from its constituent axioms --- the \(\mathrm{AB5}\) exactness
of filtered colimits and the existence of a separator.  This route is heavier
and is invoked only by \(R^if_*\).

The blocks below first record the relevant Mathlib results as dependency
anchors, then assemble the project-side interface: the light abelian-group-sheaf
bridge, the heavy module-category bridge, the derived pushforward and sheaf
cohomology, the base-change comparison morphism, the stalk-local isomorphism
criterion, and the exterior/symmetric power and local-freeness machinery.

\subsection{Mathlib substrate anchors}

\begin{definition}[Category of \(\cO_X\)-modules]
  \label{mathlib:oxmod_category}
  \lean{AlgebraicGeometry.Scheme.Modules}
  \mathlibok
  \textit{Provided by Mathlib.}
  For a scheme \(X\), Mathlib provides the category \(X.\mathrm{Modules}\) of
  sheaves of \(\cO_X\)-modules, realised as \(\mathrm{SheafOfModules}\,
  X.\mathrm{ringCatSheaf}\): sheaves of modules over the structure ring sheaf of
  \(X\) on the site of opens.  It is an abelian category, so kernels, cokernels,
  images and short exact sequences are available, and additive functors out of
  it may be right-derived.
\end{definition}

\begin{definition}[Additive pushforward of \(\cO\)-modules]
  \label{mathlib:oxmod_pushforward}
  \lean{AlgebraicGeometry.Scheme.Modules.pushforward}
  \mathlibok
  \textit{Provided by Mathlib.}
  For a morphism of schemes \(f:X\to Y\), Mathlib provides the additive
  pushforward functor
  \[
    f_*:X.\mathrm{Modules}\to Y.\mathrm{Modules},
  \]
  the degree-zero direct image of a sheaf of \(\cO_X\)-modules.  Its
  right-derived functors are the higher direct images.
\end{definition}

\begin{definition}[Right-derived functor of an additive functor]
  \label{mathlib:right_derived}
  \lean{CategoryTheory.Functor.rightDerived}
  \mathlibok
  \textit{Provided by Mathlib.}
  Let \(F:\mathcal A\to\mathcal B\) be an additive functor between abelian
  categories, with \(\mathcal A\) having enough injectives.  Mathlib provides the
  \(i\)-th right-derived functor \(F.\mathrm{rightDerived}\,i:\mathcal A\to
  \mathcal B\), computed by applying \(F\) to an injective resolution and taking
  the \(i\)-th cohomology.  A short exact sequence in \(\mathcal A\) induces the
  associated long exact sequence of the \(F.\mathrm{rightDerived}\,i\).
\end{definition}

\begin{definition}[Exterior power functor on a module category]
  \label{mathlib:exterior_power_functor}
  \lean{ModuleCat.exteriorPower.functor}
  \mathlibok
  \textit{Provided by Mathlib.}
  For a commutative ring \(R\) and \(i\in\mathbb N\), Mathlib provides the
  \(i\)-th exterior power as a functor
  \(\wedge^i:\mathrm{Mod}_R\to\mathrm{Mod}_R\) on the category of \(R\)-modules,
  natural in the module argument.  Note there is no analogous symmetric-power
  functor in Mathlib (see \cref{found:modulecat_sym_power}).
\end{definition}

\begin{definition}[Short exact sequence in an abelian category]
  \label{mathlib:short_exact}
  \lean{CategoryTheory.ShortComplex.ShortExact}
  \mathlibok
  \textit{Provided by Mathlib.}
  Mathlib encodes a short exact sequence in an abelian category as the predicate
  \(\mathrm{ShortExact}\) on a short complex \(A\xrightarrow{f}B\xrightarrow{g}C\):
  \(f\) is a monomorphism, \(g\) is an epimorphism, and the complex is exact at
  \(B\).  This is the formalism in which the kernel-bundle and power-filtration
  sequences of the secant-bundle construction are stated.
\end{definition}

\begin{lemma}[Stalkwise criterion for isomorphisms of sheaves]
  \label{mathlib:stalk_locality}
  \lean{TopCat.Presheaf.isIso_iff_stalkFunctor_map_iso}
  \mathlibok
  \textit{Provided by Mathlib.}
  A morphism \(\varphi:\mathcal F\to\mathcal G\) of sheaves on a topological
  space is an isomorphism if and only if the induced morphism on every stalk
  \(\varphi_x:\mathcal F_x\to\mathcal G_x\) is an isomorphism.
\end{lemma}

\begin{lemma}[A Grothendieck abelian category has enough injectives]
  \label{mathlib:groth_enough_injectives}
  \lean{CategoryTheory.IsGrothendieckAbelian.enoughInjectives}
  \mathlibok
  \textit{Provided by Mathlib.}
  Every Grothendieck abelian category has enough injectives: each object admits a
  monomorphism into an injective object.
\end{lemma}

\begin{lemma}[Enough injectives gives injective resolutions]
  \label{mathlib:enough_inj_has_resolutions}
  \lean{CategoryTheory.InjectiveResolution.instHasInjectiveResolution}
  \mathlibok
  \textit{Provided by Mathlib.}
  In an abelian category with enough injectives, every object admits an injective
  resolution; i.e. the category satisfies \(\mathrm{HasInjectiveResolutions}\).
\end{lemma}

\begin{lemma}[Sheaves valued in a Grothendieck category are Grothendieck abelian]
  \label{mathlib:sheaf_groth_abelian}
  \lean{CategoryTheory.Sheaf.isGrothendieckAbelian_of_essentiallySmall}
  \mathlibok
  \textit{Provided by Mathlib.}
  For an essentially small category \(\mathcal C\) equipped with a Grothendieck
  topology \(J\), and a Grothendieck abelian target category \(\mathcal A\)
  admitting sheafification (\(\mathrm{HasSheafify}\)), the category
  \(\mathrm{Sheaf}\,J\,\mathcal A\) of \(\mathcal A\)-valued sheaves on
  \(\mathcal C\) is again Grothendieck abelian.  Its valid application target in
  this chapter is the category of \emph{abelian-group}-valued sheaves on the
  essentially small site of opens of a scheme (\cref{found:absheaf_groth_abelian}),
  which literally has this form with \(\mathcal A=\mathrm{Ab}\).  It does
  \emph{not} apply to \(X.\mathrm{Modules}\): sheaves of modules over a varying
  ring sheaf are not objects of any single \(\mathrm{Sheaf}\,J\,\mathcal A\), so
  the module-category bridge (\cref{found:modules_grothendieck_abelian}) takes a
  different route.
\end{lemma}

\begin{remark}[An alternative \(\mathrm{AB5}\) input from the module category]
  For a ring \(R\), the category \(\mathrm{Mod}_R\) itself satisfies
  Grothendieck's \(\mathrm{AB5}\) axiom: it has filtered colimits and these are
  exact (they preserve finite limits / short exact sequences).  One could in
  principle feed \(\mathrm{AB5}\) for (pre)sheaves of modules from this module-level
  axiom, transported objectwise and then through sheafification.  The project does
  \emph{not} take this route: its presheaf-of-modules \(\mathrm{AB5}\) is obtained
  instead from the \(\mathrm{Ab}\)-valued functor category, via
  \cref{mathlib:addcommgrp_ab5}.
\end{remark}

\begin{lemma}[The category of abelian groups satisfies \(\mathrm{AB5}\)]
  \label{mathlib:addcommgrp_ab5}
  \lean{instAB5AddCommGrpCat}
  \mathlibok
  \textit{Provided by Mathlib.}
  The category \(\mathrm{Ab}\) of abelian groups satisfies Grothendieck's
  \(\mathrm{AB5}\) axiom: it has filtered colimits and these are exact.  This is
  the concrete base-category input from which \(\mathrm{AB5}\) for presheaves of
  modules is obtained, by lifting to the functor category \(\mathcal C^{\mathrm{op}}
  \to\mathrm{Ab}\) and transporting along the forgetful functor.
\end{lemma}

\begin{lemma}[Sheafification of presheaves of modules is exact]
  \label{mathlib:sheafification_modules_exact}
  \lean{PresheafOfModules.sheafification}
  \mathlibok
  \textit{Provided by Mathlib.}
  Mathlib provides the sheafification functor
  \[
    (-)^{\sharp}:\mathrm{PresheafOfModules}\,R\to\mathrm{SheafOfModules}\,R ,
  \]
  left adjoint to the forgetful inclusion.  It is exact: being a left adjoint it
  preserves all colimits, and Mathlib furthermore proves it preserves finite
  limits, so it is left exact as well.  This exactness is what lets one transport \(\mathrm{AB5}\) and
  the exterior/symmetric powers from the presheaf level to the sheaf level.
\end{lemma}

\begin{remark}[Why the global free sheaves do not separate]
  For a sheaf of rings \(R\) on a site and an index type \(I\) there is a free
  sheaf of modules \(R^{(I)}\), together with the natural bijection
  \(\mathrm{Hom}(R^{(I)},\mathcal M)\cong (I\to\mathcal M.\mathrm{sections})\)
  identifying morphisms out of a free sheaf with families of \emph{global}
  sections.  Because this universal property sees only global sections, the family
  of these global free sheaves is \emph{not} separating on a general scheme:
  separation by \(\{R^{(I)}\}\) would force the global-sections functor to be
  conservative, which fails already for \(\cO(-1)\) on \(\mathbb P^1\).  The
  separator route for sheaves of modules therefore uses a different family, one
  that sees sections over every open rather than only the global ones --- namely
  the free-Yoneda family \cref{mathlib:presheafmod_freeYoneda}.
\end{remark}

\begin{lemma}[Free-Yoneda separating family of presheaves of modules]
  \label{mathlib:presheafmod_freeYoneda}
  \lean{PresheafOfModules.freeYoneda.isSeparating}
  \mathlibok
  \textit{Provided by Mathlib.}
  For a presheaf of rings \(R\) on a category \(\mathcal C\), Mathlib provides the
  family \(\mathrm{freeYoneda}\,R\) of free presheaves of modules built on the
  representable (Yoneda) presheaves: for an open \(U\), a morphism
  \(\mathrm{freeYoneda}(U)\to\mathcal F\) is precisely a section of \(\mathcal F\)
  over \(U\).  This family therefore sees sections over \emph{every} open, and
  Mathlib proves it is a small separating family:
  \(\mathrm{freeYoneda}\,R\) is \(\mathrm{IsSeparating}\) (and \(\mathrm{Small}\))
  in \(\mathrm{PresheafOfModules}\,R\).
\end{lemma}

\begin{lemma}[Sheafification adjunction for presheaves of modules]
  \label{mathlib:sheafification_adjunction}
  \lean{PresheafOfModules.sheafificationAdjunction}
  \mathlibok
  \textit{Provided by Mathlib.}
  For a sheaf of rings \(R\) on a site, Mathlib provides the sheafification adjunction
  \[
    (-)^{\sharp}\dashv \iota:\;\mathrm{PresheafOfModules}\,R.\mathrm{obj}
    \;\rightleftarrows\;\mathrm{SheafOfModules}\,R ,
  \]
  whose left adjoint is sheafification and whose right adjoint \(\iota\) is the
  (faithful) forgetful inclusion of sheaves of modules into presheaves of modules.
\end{lemma}

\begin{lemma}[A coproduct of a separating family is a separator]
  \label{mathlib:isSeparator_sigma}
  \lean{CategoryTheory.isSeparator_sigma}
  \mathlibok
  \textit{Provided by Mathlib.}
  In a category with the relevant coproduct, a family of objects
  \((S_b)_{b\in\beta}\) is separating if and only if its coproduct
  \(\coprod_{b}S_b\) is a separator.  In particular a small separating family
  yields a single separator object.
\end{lemma}

\subsection{Light cohomology bridge: abelian-group sheaves}

\begin{definition}\leanok
[Underlying abelian-group sheaf of an \(\cO_X\)-module]
  \label{found:modules_to_absheaf}
  \lean{AlgebraicGeometry.Scheme.Modules.toAbSheaf}
  \uses{mathlib:oxmod_category}
  For a scheme \(X\), the forgetful functor
  \[
    (-)^{\mathrm{Ab}}:X.\mathrm{Modules}\to\mathrm{Sh}(X,\mathrm{Ab})
  \]
  sending an \(\cO_X\)-module \(\mathcal F\) to its underlying sheaf of abelian
  groups, an object of the category \(\mathrm{Sh}(X,\mathrm{Ab})\).  This is the
  forgetful functor that retains the sheaf condition and the abelian-group
  structure of \(\mathcal F\) while forgetting the \(\cO_X\)-module structure.
\end{definition}

\begin{lemma}\leanok
[Abelian-group sheaves on a scheme are Grothendieck abelian]
  \label{found:absheaf_groth_abelian}
  \lean{AlgebraicGeometry.instIsGrothendieckAbelianTopCatSheafAb}
  \uses{mathlib:sheaf_groth_abelian}
  For a scheme \(X\), the category \(\mathrm{Sh}(X,\mathrm{Ab})\) of
  abelian-group-valued sheaves on the site of opens of \(X\) is Grothendieck
  abelian.
\end{lemma}

\begin{proof}
  \uses{mathlib:sheaf_groth_abelian}
  \leanok
  The category \(\mathrm{Ab}\) of abelian groups is Grothendieck abelian, the
  site of opens of \(X\) is essentially small, and sheafification of
  abelian-group presheaves exists (\(\mathrm{HasSheafify}\)).  The category
  \(\mathrm{Sh}(X,\mathrm{Ab})\) is exactly \(\mathrm{Sheaf}\,J\,\mathrm{Ab}\) for
  the canonical Grothendieck topology \(J\) on the opens of \(X\); applying
  \cref{mathlib:sheaf_groth_abelian} to this essentially small site with target
  \(\mathcal A=\mathrm{Ab}\) yields the Grothendieck-abelian structure.
\end{proof}

\begin{lemma}\leanok
[Abelian-group sheaves have injective resolutions]
  \label{found:absheaf_has_injective_resolutions}
  \lean{AlgebraicGeometry.instHasInjectiveResolutionsTopCatSheafAb}
  \uses{found:absheaf_groth_abelian, found:absheaf_enough_injectives, mathlib:groth_enough_injectives, mathlib:enough_inj_has_resolutions}
  For a scheme \(X\), the abelian category \(\mathrm{Sh}(X,\mathrm{Ab})\) has
  injective resolutions: every abelian-group sheaf embeds into an injective one,
  and admits an injective resolution.
\end{lemma}

\begin{proof}
  \uses{found:absheaf_groth_abelian, found:absheaf_enough_injectives, mathlib:groth_enough_injectives, mathlib:enough_inj_has_resolutions}
  \leanok
  By \cref{found:absheaf_groth_abelian} the category \(\mathrm{Sh}(X,\mathrm{Ab})\)
  is Grothendieck abelian, hence has enough injectives by
  \cref{mathlib:groth_enough_injectives}; an abelian category with enough
  injectives has injective resolutions by
  \cref{mathlib:enough_inj_has_resolutions}.
\end{proof}

\begin{lemma}\leanok
[Abelian-group sheaves have enough injectives]
  \label{found:absheaf_enough_injectives}
  \lean{AlgebraicGeometry.instEnoughInjectivesTopCatSheafAb}
  \uses{found:absheaf_groth_abelian, mathlib:groth_enough_injectives}
  For a scheme \(X\), the category \(\mathrm{Sh}(X,\mathrm{Ab})\) of
  abelian-group sheaves has enough injectives: every abelian-group sheaf admits a
  monomorphism into an injective one.
\end{lemma}

\begin{proof}
  \uses{found:absheaf_groth_abelian, mathlib:groth_enough_injectives}
  \leanok
  Being Grothendieck abelian (\cref{found:absheaf_groth_abelian}), the category
  \(\mathrm{Sh}(X,\mathrm{Ab})\) has enough injectives by
  \cref{mathlib:groth_enough_injectives}.
\end{proof}

\begin{definition}\leanok
[Sheaf-level global-sections functor]
  \label{found:absheaf_global_sections}
  \lean{AlgebraicGeometry.Scheme.absheafGlobalSections}
  \uses{found:modules_to_absheaf}
  For a scheme \(X\), the sheaf-level global-sections functor
  \[
    \Gamma:\mathrm{Sh}(X,\mathrm{Ab})\to\mathrm{Ab}
  \]
  on abelian-group sheaves, given by forgetting an abelian-group sheaf to its
  underlying presheaf and evaluating at the terminal open \(\top=X\).  This is the
  functor whose right-derived functors compute sheaf cohomology
  (\cref{def:peer_scheme_hmodule}): the cohomology of an \(\cO_X\)-module is
  obtained by first passing to its underlying abelian-group sheaf
  (\cref{found:modules_to_absheaf}) and then right-deriving this \(\Gamma\).  It
  is distinct from the module-level global-sections functor of
  \cref{found:global_sections_ab}, which lands in \(\mathrm{Ab}\) directly from
  \(X.\mathrm{Modules}\).
\end{definition}

\begin{lemma}\leanok
[Sheaf-level global-sections functor is additive]
  \label{found:absheaf_global_sections_additive}
  \lean{AlgebraicGeometry.Scheme.absheafGlobalSections_additive}
  \uses{found:absheaf_global_sections}
  For a scheme \(X\), the sheaf-level global-sections functor \(\Gamma\) of
  \cref{found:absheaf_global_sections} is additive, hence may be right-derived.
\end{lemma}

\begin{proof}
  \uses{found:absheaf_global_sections}
  \leanok
  \(\Gamma\) is the composite of the additive forgetful functor from
  abelian-group sheaves to abelian-group presheaves with the additive evaluation
  functor at the open \(\top=X\); a composite of additive functors is additive.
\end{proof}

% NOTE: This module-level \(\Gamma_{\mathrm{Ab}}\) is distinct from the
% sheaf-level \(\Gamma\) (\cref{found:absheaf_global_sections}) that
% \cref{def:peer_scheme_hmodule} right-derives.  Whether right-deriving
% \(\Gamma_{\mathrm{Ab}}:X.\mathrm{Modules}\to\mathrm{Ab}\) directly recovers
% the same cohomology groups as the sheaf-level route is not yet established.
\begin{definition}[Additive global-sections functor]
  \leanok
  \label{found:global_sections_ab}
  \lean{AlgebraicGeometry.Scheme.globalSectionsAb}
  \uses{mathlib:oxmod_category}
  For a scheme \(X\), the additive global-sections functor
  \[
    \Gamma_{\mathrm{Ab}}:X.\mathrm{Modules}\to\mathrm{Ab}
  \]
  sends an \(\cO_X\)-module \(\mathcal F\) to its abelian group of global
  sections \(\mathcal F(X)\) and a morphism to its induced map on global
  sections.  It factors through the underlying-abelian-group-sheaf functor of
  \cref{found:modules_to_absheaf} followed by the global-sections functor
  \(\Gamma(X,-)\).  This module-level \(\Gamma_{\mathrm{Ab}}:X.\mathrm{Modules}\to
  \mathrm{Ab}\) is \emph{not} the functor that sheaf cohomology right-derives:
  \(H^i\) is computed by right-deriving the sheaf-level global sections
  \(\Gamma:\mathrm{Sh}(X,\mathrm{Ab})\to\mathrm{Ab}\)
  (\cref{found:absheaf_global_sections}) after passing to the underlying
  abelian-group sheaf.  Whether right-deriving \(\Gamma_{\mathrm{Ab}}\)
  directly produces the same cohomology groups as the sheaf-level route of
  \cref{def:peer_scheme_hmodule} is an open comparison question.
\end{definition}

\subsection{Heavy bridge: module category is Grothendieck abelian}

\begin{lemma}\leanok
[Presheaves of modules satisfy \(\mathrm{AB5}\)]
  \label{found:presheafmod_ab5}
  \lean{AlgebraicGeometry.instAB5OfSizePresheafOfModules}
  \uses{mathlib:addcommgrp_ab5}
  For a presheaf of rings \(R\) on the site of opens of a scheme, the category
  \(\mathrm{PresheafOfModules}\,R\) satisfies \(\mathrm{AB5OfSize}\): it has
  filtered colimits and these are exact.
\end{lemma}

\begin{proof}
  \uses{mathlib:addcommgrp_ab5}
  \leanok
  The exactness of filtered colimits is transported from \(\mathrm{Ab}\) along the
  forgetful functor to abelian-group presheaves.  The category \(\mathrm{Ab}\)
  satisfies \(\mathrm{AB5}\) (\cref{mathlib:addcommgrp_ab5}); exactness of filtered
  colimits is inherited by any functor category, so the presheaf category
  \(\mathcal C^{\mathrm{op}}\to\mathrm{Ab}\) of abelian-group presheaves again
  satisfies \(\mathrm{AB5}\).  The forgetful functor
  \(\mathrm{PresheafOfModules}\,R\to(\mathcal C^{\mathrm{op}}\to\mathrm{Ab})\)
  preserves colimits while reflecting and preserving finite limits, so the
  domain-transport principle for exact colimits (\(\mathrm{HasExactColimitsOfShape}
  \) along a functor) carries the \(\mathrm{AB5}\) exactness back to
  \(\mathrm{PresheafOfModules}\,R\).
\end{proof}

\begin{lemma}\leanok
[Sheaves of modules satisfy \(\mathrm{AB5}\)]
  \label{found:sheafmod_ab5}
  \lean{AlgebraicGeometry.instAB5OfSizeSheafOfModules}
  \uses{found:presheafmod_ab5, mathlib:sheafification_modules_exact}
  For a sheaf of rings \(R\) on the site of opens of a scheme, the category
  \(\mathrm{SheafOfModules}\,R\) satisfies \(\mathrm{AB5OfSize}\).
\end{lemma}

\begin{proof}
  \uses{found:presheafmod_ab5, mathlib:sheafification_modules_exact}
  \leanok
  Filtered colimits in \(\mathrm{SheafOfModules}\,R\) are computed by sheafifying
  the corresponding filtered colimit of presheaves.  By
  \cref{found:presheafmod_ab5} filtered colimits are exact at the presheaf level,
  and by \cref{mathlib:sheafification_modules_exact} the sheafification functor is
  exact; composing an exact filtered-colimit functor with the exact
  sheafification (and using that the forgetful inclusion is left exact) transports
  the \(\mathrm{AB5}\) exactness from presheaves to sheaves.
\end{proof}

\begin{lemma}\leanok
[Separating families transport along a left adjoint with faithful right adjoint]
  \label{found:separating_transport_left_adjoint}
  \lean{AlgebraicGeometry.isSeparating_ofObj_leftAdjoint}
  Let \(F\dashv G\) be an adjunction with \(G\) faithful.  If \(\{S_b\}_{b\in\beta}\)
  is a separating family in the source category, then \(\{F(S_b)\}_{b\in\beta}\)
  is a separating family in the target category.
\end{lemma}

\begin{proof}

\leanok
  Proved directly in Lean.
\end{proof}

\begin{lemma}\leanok
[Sheaves of modules have a separator]
  \label{found:sheafmod_separator}
  \lean{AlgebraicGeometry.instHasSeparatorSheafOfModules}
  \uses{mathlib:presheafmod_freeYoneda, mathlib:sheafification_adjunction, mathlib:isSeparator_sigma, found:separating_transport_left_adjoint}
  For a sheaf of rings \(R\) on the site of opens of a scheme, the category
  \(\mathrm{SheafOfModules}\,R\) has a separator.
\end{lemma}

\begin{proof}
  \uses{mathlib:presheafmod_freeYoneda, mathlib:sheafification_adjunction, mathlib:isSeparator_sigma, found:separating_transport_left_adjoint}
  \leanok
  Work at the presheaf level first.  The free-Yoneda family
  \(\mathrm{freeYoneda}\,R.\mathrm{obj}\) of free presheaves of modules on the
  representable presheaves is a small separating family of
  \(\mathrm{PresheafOfModules}\,R.\mathrm{obj}\)
  (\cref{mathlib:presheafmod_freeYoneda}): a morphism out of
  \(\mathrm{freeYoneda}(U)\) is exactly a section over the open \(U\), so the
  family detects morphisms by their effect on sections over \emph{every} open, not
  merely the global ones.

  Transport this family across the sheafification adjunction
  \((-)^{\sharp}\dashv\iota\) of \cref{mathlib:sheafification_adjunction}.  A separating family transports
  along a left adjoint whose right adjoint is faithful
  (\cref{found:separating_transport_left_adjoint}): since \(\iota\) is faithful,
  the images \(\bigl(\mathrm{freeYoneda}(U)\bigr)^{\sharp}\) form a separating
  family of \(\mathrm{SheafOfModules}\,R\).

  Finally, the coproduct of a separating family is a separator
  (\cref{mathlib:isSeparator_sigma}), so the coproduct of the sheafified
  free-Yoneda family is a single separator object of
  \(\mathrm{SheafOfModules}\,R\).
\end{proof}

\begin{lemma}

\leanok
[Sheaves of modules over a sheaf of rings form a Grothendieck abelian category]
  \label{found:sheafmod_grothendieck_abelian}
  \lean{AlgebraicGeometry.instIsGrothendieckAbelianSheafOfModules}
  \uses{found:sheafmod_ab5, found:sheafmod_separator}
  For a sheaf of rings \(R\) on the (small) site of opens of a scheme, the category
  \(\mathrm{SheafOfModules}\,R\) is Grothendieck abelian.
\end{lemma}

\begin{proof}
  \uses{found:sheafmod_ab5, found:sheafmod_separator}
  \leanok
  Being Grothendieck abelian amounts to four data on the abelian category
  \(\mathrm{SheafOfModules}\,R\): existence of filtered colimits, local smallness,
  exactness of filtered colimits (\(\mathrm{AB5OfSize}\)), and the existence of a
  separator.  The first two hold by inference (filtered colimits are computed via
  sheafification of objectwise colimits, and the category is locally small).  The
  substantive inputs are the \(\mathrm{AB5}\) exactness, supplied by
  \cref{found:sheafmod_ab5}, and the separator, supplied by
  \cref{found:sheafmod_separator}.  Assembling the four fields yields the
  Grothendieck-abelian structure.  The essentially-small theorem
  \cref{mathlib:sheaf_groth_abelian} is \emph{not} used, as sheaves of modules over
  a varying ring sheaf are not of the form \(\mathrm{Sheaf}\,J\,\mathcal A\).
\end{proof}

\begin{lemma}\leanok
[\(\cO_X\)-modules form a Grothendieck abelian category]
  \label{found:modules_grothendieck_abelian}
  \lean{AlgebraicGeometry.instIsGrothendieckAbelianModules}
  \uses{found:sheafmod_grothendieck_abelian, mathlib:oxmod_category}
  For every scheme \(X\), the category \(X.\mathrm{Modules}\) is Grothendieck
  abelian.
\end{lemma}

\begin{proof}
  \uses{found:sheafmod_grothendieck_abelian, mathlib:oxmod_category}
  \leanok
  By definition \(X.\mathrm{Modules}=\mathrm{SheafOfModules}\,X.\mathrm{ringCatSheaf}\)
  (\cref{mathlib:oxmod_category}), so this is the specialisation of the general
  instance \cref{found:sheafmod_grothendieck_abelian} to the structure ring sheaf
  \(R=X.\mathrm{ringCatSheaf}\), transferred across that definitional equality.  The
  four-field assembly of the Grothendieck-abelian structure (filtered colimits and
  local smallness by inference, \(\mathrm{AB5}\), and the separator) lives in the
  general instance.
\end{proof}

\begin{lemma}\leanok
[\(\cO_X\)-modules have injective resolutions]
  \label{found:has_injective_resolutions}
  \lean{AlgebraicGeometry.instHasInjectiveResolutionsModules}
  \uses{found:modules_grothendieck_abelian, mathlib:groth_enough_injectives, mathlib:enough_inj_has_resolutions}
  For every scheme \(X\), the abelian category \(X.\mathrm{Modules}\) has enough
  injectives, so every \(\cO_X\)-module admits an injective resolution; in
  particular, the hypothesis of \cref{mathlib:right_derived} is satisfied for
  the pushforward functor in the module category.
\end{lemma}

\begin{proof}
  \uses{found:modules_grothendieck_abelian, mathlib:groth_enough_injectives, mathlib:enough_inj_has_resolutions}
  \leanok
  By \cref{found:modules_grothendieck_abelian} the category \(X.\mathrm{Modules}\)
  is Grothendieck abelian, hence has enough injectives by
  \cref{mathlib:groth_enough_injectives}; an abelian category with enough
  injectives has injective resolutions by
  \cref{mathlib:enough_inj_has_resolutions}.
\end{proof}

\subsection{Exterior and symmetric powers}

\begin{lemma}\leanok
[Commutative-ring structure on the structure-sheaf sections]
  \label{found:ringcatsheaf_commring}
  \lean{AlgebraicGeometry.Scheme.ringCatSheaf_commRing}
  \uses{mathlib:oxmod_category}
  For a scheme \(X\) and an open \(U\), the value \(X.\mathrm{ringCatSheaf}(U)\) of
  the structure ring sheaf --- a priori only an object of \(\mathrm{RingCat}\) ---
  carries a commutative-ring structure, compatible with (indeed identical to) the
  commutative-ring structure of the structure-sheaf sections \(\cO_X(U)=
  X.\mathrm{presheaf}(U)\) as an object of \(\mathrm{CommRingCat}\).
\end{lemma}

\begin{proof}
  \uses{mathlib:oxmod_category}
  \leanok
  The ring sheaf \(X.\mathrm{ringCatSheaf}\) is the image of the structure sheaf
  \(\cO_X\) under the forgetful functor \(\mathrm{CommRingCat}\to\mathrm{RingCat}\):
  its value at \(U\) is the same underlying set with the same addition and
  multiplication as \(\cO_X(U)\).  The commutativity carried by the
  \(\mathrm{CommRingCat}\) object \(\cO_X(U)\) is therefore present on the
  sections, and is recovered as a \(\mathrm{CommRing}\) instance on the
  \(\mathrm{RingCat}\)-valued \(X.\mathrm{ringCatSheaf}(U)\), which the
  \(\mathrm{RingCat}\) type level forgets.
\end{proof}

\begin{lemma}\leanok
[Linear-algebra base change of exterior powers]
  \label{found:exterior_base_change_linear}
  \lean{ModuleCat.exteriorPower.baseChange}
  \uses{found:exterior_base_change_alt, mathlib:exterior_power_functor}
  Let \(A\) be a commutative ring, \(B\) an \(A\)-algebra, and \(N\) a
  \(B\)-module equipped with a compatible \(A\)-module structure and the scalar
  tower \(\mathrm{IsScalarTower}\,A\,B\,N\).  Then for each \(i\in\mathbb N\)
  there is an \(A\)-linear comparison map
  \[
    \textstyle\bigwedge^i_A N\to\bigwedge^i_B N
  \]
  from the exterior power formed over \(A\) to the exterior power formed over
  \(B\); on a decomposable wedge it acts as the identity on representatives:
  \(n_1\wedge\cdots\wedge n_i\mapsto n_1\wedge\cdots\wedge n_i\).  The target
  exterior power inherits the \(A\)-module structure via the scalar tower
  \(\mathrm{IsScalarTower}\,A\,B\,(\bigwedge^i_B N)\), so that source and target
  are both \(A\)-modules.
\end{lemma}

\begin{proof}
  \uses{found:exterior_base_change_alt, mathlib:exterior_power_functor}
  \leanok
  Over the fixed ring \(B\) the exterior power \(\bigwedge^i_B N\) is the object
  of \cref{mathlib:exterior_power_functor}; the new content is the
  comparison with the exterior power formed over the subring of scalars \(A\).
  The \(i\)-fold wedge \(n_1\wedge\cdots\wedge n_i\) over \(A\) is sent to the
  same expression read over \(B\); by the alternating-multilinear universal
  property of \(\bigwedge^i_A N\) this rule extends uniquely to an \(A\)-linear
  map into \(\bigwedge^i_B N\), the latter regarded as an \(A\)-module via the
  scalar tower.  The action-on-generators identity is then the defining equation
  of the map on the spanning decomposable wedges.
\end{proof}

\begin{lemma}\leanok
[Base-change functoriality of exterior powers]
  \label{found:exterior_power_base_change}
  \lean{ModuleCat.exteriorPower.baseChangeMap}
  \uses{found:exterior_base_change_linear, mathlib:exterior_power_functor}
  Let \(\varphi:R\to S\) be a homomorphism of commutative rings, \(M\) an
  \(R\)-module, \(N\) an \(S\)-module, and \(g:M\to N\) a \(\varphi\)-semilinear
  map --- equivalently an \(R\)-linear map \(M\to N{\restriction}_R\) into the
  restriction of scalars of \(N\) along \(\varphi\).  Then for each
  \(i\in\mathbb N\) there is an induced \(\varphi\)-semilinear map
  \[
    \wedge^i g:\textstyle\bigwedge^i_R M\to\bigwedge^i_S N
  \]
  on exterior powers, compatible with restriction of scalars along \(\varphi\) and
  determined on a decomposable wedge by \(n_1\wedge\cdots\wedge n_i\mapsto
  g(n_1)\wedge\cdots\wedge g(n_i)\) read over \(S\).  Its behaviour under the
  identity ring map and under composition of semilinear maps --- the naturality
  that makes base change functorial in the ring map --- is recorded separately in
  \cref{found:exterior_base_change_map_id} and
  \cref{found:exterior_base_change_map_comp}.
  This is the ingredient that Mathlib lacks: it provides the exterior power as a
  functor over a \emph{fixed} commutative base
  (\cref{mathlib:exterior_power_functor}) but not its variance in the base ring,
  so this base-change functoriality is supplied here as project infrastructure.
  It arises as the semilinear wrapper of the linear-algebra comparison map
  \cref{found:exterior_base_change_linear}, via restriction of scalars along \(\varphi\).
\end{lemma}

\begin{proof}
  \uses{found:exterior_base_change_linear, mathlib:exterior_power_functor}
  \leanok
  A \(\varphi\)-semilinear map \(g:M\to N\) is an \(R\)-linear map
  \(M\to N{\restriction}_R\), so applying the fixed-base functor
  \cref{mathlib:exterior_power_functor} over \(R\) yields an \(R\)-linear map
  \(\bigwedge^i_R M\to\bigwedge^i_R(N{\restriction}_R)\).  It remains to compare
  \(\bigwedge^i_R(N{\restriction}_R)\) with \(\bigl(\bigwedge^i_S N\bigr){\restriction}_R\);
  this is exactly the linear-algebra comparison of
  \cref{found:exterior_base_change_linear}, applied with the algebra structure
  on \(S\) induced by \(\varphi\), the \(R\)-module structure on \(N\) supplied
  by restriction of scalars, and the scalar tower \(\mathrm{IsScalarTower}\,R\,S\,N\).
  Composing the two \(R\)-linear maps sends a decomposable wedge
  \(n_1\wedge\cdots\wedge n_i\) to the same decomposable wedge formed over \(S\)
  and produces the desired \(\varphi\)-semilinear map, with the stated action on
  generators.  This construction supplies only the semilinear map and its value on
  generators; that base change along the identity ring map recovers the fixed-base
  exterior power and that it respects composition of semilinear maps is established
  separately in \cref{found:exterior_base_change_map_id} and
  \cref{found:exterior_base_change_map_comp}.
\end{proof}

\begin{lemma}[Base change along the identity ring map --- exterior]\leanok
  \label{found:exterior_base_change_map_id}
  \lean{ModuleCat.exteriorPower.baseChangeMap_id}
  \uses{found:exterior_power_base_change, mathlib:exterior_power_functor}
  Let \(R\) be a commutative ring, \(M\) an \(R\)-module and \(g:M\to M\) an
  \(R\)-linear map, viewed as an \(\mathrm{id}_R\)-semilinear map.  Then the
  base-change map \(\wedge^i g\) of \cref{found:exterior_power_base_change} formed
  along the identity ring homomorphism \(\mathrm{id}_R\) coincides with the
  fixed-base exterior power \(\wedge^i g\) of the functor
  \cref{mathlib:exterior_power_functor}, both being \(R\)-linear endomorphisms of
  \(\bigwedge^i_R M\).
\end{lemma}

\begin{proof}\leanok
  \uses{found:exterior_power_base_change, mathlib:exterior_power_functor}
  Both maps are \(R\)-linear endomorphisms of \(\bigwedge^i_R M\), so it suffices
  to compare them on the decomposable wedges \(n_1\wedge\cdots\wedge n_i\), which
  span.  Base change along \(\mathrm{id}_R\) leaves the base ring unchanged, so its
  action \(n_1\wedge\cdots\wedge n_i\mapsto g(n_1)\wedge\cdots\wedge g(n_i)\) is
  exactly that of the fixed-base functor \cref{mathlib:exterior_power_functor}.
\end{proof}

\begin{lemma}[Base change respects composition of ring maps --- exterior]\leanok
  \label{found:exterior_base_change_map_comp}
  \lean{ModuleCat.exteriorPower.baseChangeMap_comp}
  \uses{found:exterior_power_base_change}
  Let \(\varphi:R\to S\) and \(\psi:S\to T\) be homomorphisms of commutative rings,
  let \(g:M\to N\) be \(\varphi\)-semilinear and \(h:N\to P\) be
  \(\psi\)-semilinear.  Then the base-change map of the composite \(h\circ g\)
  (which is \((\psi\circ\varphi)\)-semilinear), formed along \(\psi\circ\varphi\),
  factors as the composite of the base-change maps of \(g\) and \(h\):
  \[
    \wedge^i(h\circ g)=\wedge^i h\circ\wedge^i g
    :\textstyle\bigwedge^i_R M\to\bigwedge^i_T P .
  \]
\end{lemma}

\begin{proof}\leanok
  \uses{found:exterior_power_base_change}
  Both sides are additive maps \(\bigwedge^i_R M\to\bigwedge^i_T P\) that agree on
  the spanning decomposable wedges: each sends \(n_1\wedge\cdots\wedge n_i\) to
  \(h(g(n_1))\wedge\cdots\wedge h(g(n_i))\) read over \(T\).  Hence they coincide.
\end{proof}

\begin{lemma}\leanok
[Alternating map underlying the exterior base-change comparison]
  \label{found:exterior_base_change_alt}
  \lean{ModuleCat.exteriorPower.baseChangeAlt}
  \uses{mathlib:exterior_power_functor}
  The \(A\)-alternating map \(N^i\to\bigwedge^i_B N\) underlying the comparison map
  \cref{found:exterior_base_change_linear}: the canonical \(B\)-alternating map
  \((n_1,\dots,n_i)\mapsto n_1\wedge\cdots\wedge n_i\) into \(\bigwedge^i_B N\),
  with its scalars restricted from \(B\) to \(A\) along the algebra structure.
\end{lemma}

\begin{proof}
  \uses{mathlib:exterior_power_functor}
  \leanok
  The fixed-base exterior power over \(B\) (\cref{mathlib:exterior_power_functor})
  carries its universal \(B\)-alternating map; restriction of scalars along
  \(A\to B\) turns it into an \(A\)-alternating map, alternating because the
  underlying additive map is unchanged.
\end{proof}

\begin{lemma}\leanok
[Generator action of the exterior base-change comparison]
  \label{found:exterior_base_change_apply_iotamulti}
  \lean{ModuleCat.exteriorPower.baseChange_apply_ιMulti}
  \uses{found:exterior_base_change_linear, found:exterior_base_change_alt}
  The comparison map \cref{found:exterior_base_change_linear} sends the
  decomposable wedge \(v_1\wedge\cdots\wedge v_i\) formed over \(A\) to the same
  wedge \(v_1\wedge\cdots\wedge v_i\) formed over \(B\).
\end{lemma}

\begin{proof}
  \uses{found:exterior_base_change_linear, found:exterior_base_change_alt}
  \leanok
  This is the defining equation of \cref{found:exterior_base_change_linear} on the
  spanning decomposable wedges, read off from the universal property of
  \(\bigwedge^i_A N\) applied to the \(A\)-alternating map
  \cref{found:exterior_base_change_alt}.
\end{proof}

\begin{lemma}\leanok
[Generator action of the exterior base-change map]
  \label{found:exterior_base_change_map_apply_iotamulti}
  \lean{ModuleCat.exteriorPower.baseChangeMap_apply_ιMulti}
  \uses{found:exterior_power_base_change, found:exterior_base_change_apply_iotamulti}
  The base-change map \cref{found:exterior_power_base_change} along
  \(\varphi:R\to S\) sends the decomposable wedge \(v_1\wedge\cdots\wedge v_i\)
  over \(R\) to \(g(v_1)\wedge\cdots\wedge g(v_i)\) over \(S\), the wedge of the
  \(g\)-images of the entries.
\end{lemma}

\begin{proof}
  \uses{found:exterior_power_base_change, found:exterior_base_change_apply_iotamulti}
  \leanok
  Unfolding \cref{found:exterior_power_base_change} as the fixed-base exterior
  power of \(g\) over \(R\) followed by the comparison map
  \cref{found:exterior_base_change_linear}, the wedge \(v_1\wedge\cdots\wedge v_i\)
  maps first to \(g(v_1)\wedge\cdots\wedge g(v_i)\) over \(R\) and then, by
  \cref{found:exterior_base_change_apply_iotamulti}, to the same wedge read over
  \(S\).
\end{proof}

\begin{lemma}[Exterior power of a presheaf of modules]\leanok
  \label{found:presheafmod_exterior_power}
  \lean{AlgebraicGeometry.PresheafOfModules.exteriorPower}
  \uses{found:ringcatsheaf_commring, found:exterior_power_base_change, found:exterior_base_change_map_id, found:exterior_base_change_map_comp, mathlib:exterior_power_functor, found:exterior_base_change_map_cat, found:presheafmod_extpowmap, found:presheafmod_extpowmap_id, found:presheafmod_extpowmap_comp}
  For \(i\in\mathbb N\), the \(i\)-th exterior power functor
  \[
    \wedge^i:\mathrm{PresheafOfModules}\,R\to\mathrm{PresheafOfModules}\,R ,
  \]
  given over each open \(U\) by the module-level exterior power
  \(\mathrm{ModuleCat.exteriorPower}\) of \(\mathcal F(U)\), and functorial in the
  restriction maps.
\end{lemma}

\begin{proof}\leanok
  \uses{found:ringcatsheaf_commring, found:exterior_power_base_change, found:exterior_base_change_map_id, found:exterior_base_change_map_comp, mathlib:exterior_power_functor, found:exterior_base_change_map_cat, found:presheafmod_extpowmap, found:presheafmod_extpowmap_id, found:presheafmod_extpowmap_comp}
  Two obstructions must be cleared.  First, \(\mathrm{PresheafOfModules}\,R\) is
  \(\mathrm{RingCat}\)-valued, so the sections ring \(R(U)\) does not present as
  commutative at the type level and the module-level exterior-power functor
  \cref{mathlib:exterior_power_functor}, which requires a commutative base, does
  not apply directly; \cref{found:ringcatsheaf_commring} restores the
  \(\mathrm{CommRing}\) structure on each \(R(U)\), so the openwise exterior power
  \(\wedge^i_{R(U)}\mathcal F(U)\) is defined.  Second, the restriction map of
  \(\mathcal F\) along an inclusion \(V\subseteq U\) is not \(R(U)\)-linear but
  \emph{semilinear} over the ring map \(R(U)\to R(V)\), so the induced map on
  exterior powers must run \(\wedge^i_{R(U)}\mathcal F(U)\to\wedge^i_{R(V)}\mathcal
  F(V)\) over \(R(U)\to R(V)\); this is exactly the base-change functoriality of
  \cref{found:exterior_power_base_change}.  Applying that lemma to each
  restriction map produces semilinear transition maps on exterior powers, which
  inherit the presheaf identity and composition laws from those of
  \(\mathcal F\) --- via the identity-ring-map and composition naturality of base
  change, \cref{found:exterior_base_change_map_id} and
  \cref{found:exterior_base_change_map_comp} --- and assemble into a presheaf of
  modules \(\wedge^i\mathcal F\), functorially in \(\mathcal F\).
\end{proof}

\begin{definition}[Exterior power of an \(\cO_X\)-module]\leanok
  \label{found:oxmod_exterior_power}
  \lean{AlgebraicGeometry.Scheme.Modules.exteriorPower}
  \uses{found:presheafmod_exterior_power, mathlib:sheafification_modules_exact, mathlib:oxmod_category}
  For \(i\in\mathbb N\), the exterior power functor
  \(\wedge^i:X.\mathrm{Modules}\to X.\mathrm{Modules}\), obtained by applying the
  presheaf-level exterior power \cref{found:presheafmod_exterior_power} to the
  underlying presheaf of an \(\cO_X\)-module and then sheafifying
  (\cref{mathlib:sheafification_modules_exact}).  Kemeny's bundles
  \(\wedge^{k+1}M_L\), \(\wedge^k\pi^*\cM\) and \(\wedge^k\Gamma\) are instances
  of this functor.
\end{definition}

\begin{definition}[Symmetric power type]
  \label{mathlib:symmetricpower_type}
  \lean{SymmetricPower}
  \mathlibok
  \textit{Provided by Mathlib.}
  For a commutative (semi)ring \(R\), an index type \(\iota\) and an
  \(R\)-module \(M\), Mathlib provides the symmetric-power \emph{type}
  \(\Sym^i M=\mathrm{SymmetricPower}\,R\,\iota\,M\) --- the quotient of the
  \(\iota\)-fold tensor power by the permutation relation --- together with its
  \(R\)-module structure and the generating family of decomposable symmetric
  tensors \(\mathrm{tprod}\), whose span is everything.  It does \emph{not}
  provide the functorial action \(\Sym^i f\) of a linear map, nor any functor
  packaging; those are built in \cref{found:modulecat_sym_power}.
\end{definition}

\begin{lemma}\leanok
[Factorwise tensor map underlying \(\Sym^i f\)]
  \label{found:sympow_map_aux}
  \lean{SymmetricPower.mapAux}
  \uses{mathlib:symmetricpower_type}
  For an \(R\)-linear map \(f:M\to N\), the auxiliary \(R\)-linear map
  \[
    \textstyle\bigotimes^i M\to\Sym^i N
  \]
  obtained by applying \(f\) factor-wise on the \(i\)-fold tensor power and then
  symmetrising --- sending a decomposable tensor to the \(\mathrm{tprod}\) of the
  images of its factors.  This is the map that descends to \(\Sym^i f\).
\end{lemma}

\begin{proof}

\leanok
  Proved directly in Lean.
\end{proof}

\begin{lemma}\leanok
[The auxiliary map respects symmetrisation]
  \label{found:sympow_map_aux_respects}
  \lean{SymmetricPower.mapAux_respects}
  \uses{found:sympow_map_aux, mathlib:symmetricpower_type}
  The factorwise map of \cref{found:sympow_map_aux} is constant on the
  symmetrisation-equivalence classes: the factor-wise action commutes with
  permutations of the tensor factors, hence it descends through the
  symmetrisation quotient that defines \(\Sym^i M\).
\end{lemma}

\begin{proof}

\leanok
  Proved directly in Lean.
\end{proof}

\begin{lemma}\leanok
[Induced map \(\Sym^i f\) on symmetric powers]
  \label{found:sympow_map}
  \lean{SymmetricPower.map}
  \uses{found:sympow_map_aux, found:sympow_map_aux_respects}
  For an \(R\)-linear map \(f:M\to N\), the induced \(R\)-linear map
  \[
    \Sym^i f:\Sym^i M\to\Sym^i N ,
  \]
  obtained by descending the factorwise map \cref{found:sympow_map_aux} through
  the symmetrisation quotient; well-definedness is
  \cref{found:sympow_map_aux_respects}.
\end{lemma}

\begin{proof}

\leanok
  Proved directly in Lean.
\end{proof}

\begin{lemma}\leanok
[\(\Sym^i f\) acts factorwise on decomposable tensors]
  \label{found:sympow_map_tprod}
  \lean{SymmetricPower.map_tprod}
  \uses{found:sympow_map, found:sympow_map_aux}
  The map \(\Sym^i f\) of \cref{found:sympow_map} sends a decomposable symmetric
  tensor \(m_1\cdots m_i\) to \(f(m_1)\cdots f(m_i)\).
\end{lemma}

\begin{proof}

\leanok
  Proved directly in Lean.
\end{proof}

\begin{lemma}\leanok
[\(\Sym^i\) preserves identities]
  \label{found:sympow_map_id}
  \lean{SymmetricPower.map_id}
  \uses{found:sympow_map, found:sympow_map_tprod}
  \(\Sym^i(\mathrm{id}_M)=\mathrm{id}_{\Sym^i M}\): two linear maps agreeing on
  the spanning decomposable tensors coincide.
\end{lemma}

\begin{proof}

\leanok
  Proved directly in Lean.
\end{proof}

\begin{lemma}\leanok
[\(\Sym^i\) preserves composition]
  \label{found:sympow_map_comp}
  \lean{SymmetricPower.map_comp}
  \uses{found:sympow_map, found:sympow_map_tprod}
  \(\Sym^i(g\circ f)=\Sym^i g\circ\Sym^i f\).
\end{lemma}

\begin{proof}

\leanok
  Proved directly in Lean.
\end{proof}

\begin{definition}\leanok
[Symmetric power object in the module category]
  \label{found:modulecat_sympow_obj}
  \lean{ModuleCat.symmetricPower}
  \uses{mathlib:symmetricpower_type}
  For an object \(M\) of \(\mathrm{Mod}_R\) and \(n\in\mathbb N\), the symmetric
  power \(\Sym^n M\) as an object of \(\mathrm{Mod}_R\), with the index type taken
  to be \(\mathrm{ULift}(\mathrm{Fin}\,n)\) so that the ring/index universe
  coupling of the symmetric-power type is satisfied while \(R\) stays in a general
  universe (mirroring the exterior-power object of
  \cref{mathlib:exterior_power_functor}).
\end{definition}

\begin{definition}\leanok
[Induced morphism on symmetric power objects]
  \label{found:modulecat_sympow_map}
  \lean{ModuleCat.symmetricPower.map}
  \uses{found:modulecat_sympow_obj, found:sympow_map}
  For a morphism \(f:M\to N\) of \(\mathrm{Mod}_R\) and \(n\in\mathbb N\), the
  induced morphism \(\Sym^n M\to\Sym^n N\) of \(\mathrm{Mod}_R\), the categorical
  wrapper of the linear map \(\Sym^n f\) of \cref{found:sympow_map}.
\end{definition}

\begin{lemma}\leanok
[Symmetric power functor on a module category]
  \label{found:modulecat_sym_power}
  \lean{ModuleCat.symmetricPower.functor}
  \uses{found:modulecat_sympow_obj, found:modulecat_sympow_map, found:sympow_map_id, found:sympow_map_comp}
  For a commutative ring \(R\) and \(i\in\mathbb N\), the \(i\)-th symmetric power
  as a functor \(\Sym^i:\mathrm{Mod}_R\to\mathrm{Mod}_R\).  Mathlib supplies only
  the symmetric-power \emph{object} \(\Sym^i M\) (the type
  \(\mathrm{SymmetricPower}\,R\,\iota\,M\) of \cref{mathlib:symmetricpower_type})
  and its generating decomposable tensors; the induced linear map on a morphism
  and the functor packaging are both absent and are constructed here in two
  layers.

  \emph{Layer 1 (linear algebra).}  For an \(R\)-linear map \(f:M\to N\) one
  forms the induced map
  \[
    \Sym^i f:\Sym^i M\to\Sym^i N
  \]
  by descending the \(i\)-fold tensor-power action of \(f\) through the
  symmetrisation quotient: \(f\) acts factor-wise on the \(i\)-fold tensor power
  by tensor-power functoriality, this action carries the permutation relation
  into itself, so it passes to the quotient and defines \(\Sym^i f\) on
  decomposable symmetric tensors by \(m_1\cdots m_i\mapsto f(m_1)\cdots f(m_i)\).
  The functoriality laws
  \[
    \Sym^i(\mathrm{id}_M)=\mathrm{id}_{\Sym^i M},\qquad
    \Sym^i(g\circ f)=\Sym^i g\circ\Sym^i f
  \]
  follow from the corresponding identities for the tensor power together with the
  fact that the decomposable symmetric tensors span \(\Sym^i M\), so two linear
  maps agreeing on them coincide.

  \emph{Layer 2 (module category).}  Wrapping Layer 1 gives the functor
  \(\Sym^i:\mathrm{Mod}_R\to\mathrm{Mod}_R\), \(M\mapsto\Sym^i M\),
  \(f\mapsto\Sym^i f\), mirroring the exterior-power functor
  \cref{mathlib:exterior_power_functor}; the functor laws are exactly the
  Layer 1 identities.
\end{lemma}

\begin{proof}
  \uses{found:modulecat_sympow_obj, found:modulecat_sympow_map, found:sympow_map_id, found:sympow_map_comp, mathlib:exterior_power_functor}
  \leanok
  Layer~1: the symmetrisation quotient defining \(\Sym^i M\) is taken on the
  \(i\)-fold tensor power, whose functorial action sends \(f:M\to N\) to the
  factor-wise map on tensors.  This map intertwines the permutation relations on
  source and target, hence descends to a well-defined linear map
  \(\Sym^i f:\Sym^i M\to\Sym^i N\) determined by
  \(m_1\cdots m_i\mapsto f(m_1)\cdots f(m_i)\) on decomposable symmetric tensors.
  Because these tensors span \(\Sym^i M\), the identities
  \(\Sym^i(\mathrm{id})=\mathrm{id}\) and
  \(\Sym^i(g\circ f)=\Sym^i g\circ\Sym^i f\) reduce to the tensor-power versions,
  which they inherit factor-wise.  Layer~2: \(M\mapsto\Sym^i M\),
  \(f\mapsto\Sym^i f\) is then a functor \(\mathrm{Mod}_R\to\mathrm{Mod}_R\), with
  the functor laws supplied by the Layer~1 identities, in the same packaging as
  \cref{mathlib:exterior_power_functor}.
\end{proof}

\begin{lemma}\leanok
[Linear-algebra base change of symmetric powers]
  \label{found:sympow_base_change_linear}
  \lean{SymmetricPower.baseChange}
  \uses{found:sympow_base_change_aux, mathlib:symmetricpower_type}
  Let \(A\) be a commutative ring, \(B\) an \(A\)-algebra, and \(N\) a
  \(B\)-module equipped with a compatible \(A\)-module structure and the scalar
  tower \(\mathrm{IsScalarTower}\,A\,B\,N\).  Then for each \(i\in\mathbb N\)
  there is an \(A\)-linear comparison map
  \[
    \Sym^i_A N\to\Sym^i_B N
  \]
  from the symmetric power formed over \(A\) to the symmetric power formed over
  \(B\); on a decomposable symmetric tensor it acts as the identity on
  representatives: \(n_1\cdots n_i\mapsto n_1\cdots n_i\).  The target symmetric
  power inherits the \(A\)-module structure via the scalar tower
  \(\mathrm{IsScalarTower}\,A\,B\,(\Sym^i_B N)\), so that source and target are
  both \(A\)-modules.
\end{lemma}

\begin{proof}
  \uses{mathlib:symmetricpower_type, found:sympow_base_change_aux, found:sympow_base_change_aux_respects}
  \leanok
  Over the fixed ring \(B\) the symmetric power \(\Sym^i_B N\) is the object of
  \cref{mathlib:symmetricpower_type}; the new content is the comparison with the
  symmetric power formed over the subring of scalars \(A\).  The decomposable
  symmetric tensor \(n_1\cdots n_i\) over \(A\) is sent to the same expression
  read over \(B\).  This map arises from the factorwise lift
  \cref{found:sympow_base_change_aux} on the \(i\)-fold tensor power, which respects
  symmetrisation (\cref{found:sympow_base_change_aux_respects}): the multilinear assignment
  \((n_1,\dots,n_i)\mapsto n_1\cdots n_i\in\Sym^i_B N\), viewed as an
  \(A\)-multilinear map via restriction of scalars, factors through \(\Sym^i_A N\)
  by the universal property of the symmetric power and the symmetrisation quotient,
  yielding the desired \(A\)-linear map.  The action-on-generators identity is the
  defining equation of the map on the spanning decomposable symmetric tensors.
\end{proof}

\begin{lemma}\leanok
[Base-change functoriality of symmetric powers]
  \label{found:sym_power_base_change}
  \lean{ModuleCat.symmetricPower.baseChangeMap}
  \uses{found:modulecat_sym_power, found:sympow_base_change_linear}
  Let \(\varphi:R\to S\) be a homomorphism of commutative rings, \(M\) an
  \(R\)-module, \(N\) an \(S\)-module, and \(g:M\to N\) a \(\varphi\)-semilinear
  map --- equivalently an \(R\)-linear map \(M\to N{\restriction}_R\) into the
  restriction of scalars of \(N\) along \(\varphi\).  Then for each
  \(i\in\mathbb N\) there is an induced \(\varphi\)-semilinear map
  \[
    \Sym^i g:\Sym^i_R M\to\Sym^i_S N
  \]
  on symmetric powers, compatible with restriction of scalars along \(\varphi\)
  and determined on a decomposable symmetric tensor by
  \(m_1\cdots m_i\mapsto g(m_1)\cdots g(m_i)\) read over \(S\).  Its behaviour
  under the identity ring map and under composition of semilinear maps --- the
  naturality that makes base change functorial in the ring map --- is recorded
  separately in \cref{found:sym_base_change_map_id} and
  \cref{found:sym_base_change_map_comp}.  Exactly as on the exterior side
  (\cref{found:exterior_power_base_change}), Mathlib supplies no variance of the
  symmetric power in the base ring, so this base-change functoriality is project
  infrastructure; it arises as the semilinear wrapper of the linear-algebra comparison
  map \cref{found:sympow_base_change_linear}, via restriction of scalars along \(\varphi\).
\end{lemma}

\begin{proof}
  \uses{found:modulecat_sym_power, found:sympow_base_change_linear}
  \leanok
  A \(\varphi\)-semilinear map \(g:M\to N\) is an \(R\)-linear map
  \(M\to N{\restriction}_R\), so applying the fixed-base symmetric-power functor
  \cref{found:modulecat_sym_power} over \(R\) yields an \(R\)-linear map
  \(\Sym^i_R M\to\Sym^i_R(N{\restriction}_R)\).  It remains to compare
  \(\Sym^i_R(N{\restriction}_R)\) with \(\bigl(\Sym^i_S N\bigr){\restriction}_R\);
  this is exactly the linear-algebra comparison of
  \cref{found:sympow_base_change_linear}, applied with the algebra structure on
  \(S\) induced by \(\varphi\), the \(R\)-module structure on \(N\) supplied by
  restriction of scalars, and the scalar tower \(\mathrm{IsScalarTower}\,R\,S\,N\).
  Composing the two \(R\)-linear maps sends a decomposable symmetric tensor
  \(m_1\cdots m_i\) to
  \(g(m_1)\cdots g(m_i)\) formed over \(S\) and produces the desired
  \(\varphi\)-semilinear map, with the stated action on generators.  This
  construction supplies only the semilinear map and its value on generators; that
  base change along the identity ring map recovers the fixed-base symmetric power
  and that it respects composition of semilinear maps is established separately in
  \cref{found:sym_base_change_map_id} and \cref{found:sym_base_change_map_comp}.
\end{proof}

\begin{lemma}[Base change along the identity ring map --- symmetric]\leanok
  \label{found:sym_base_change_map_id}
  \lean{ModuleCat.symmetricPower.baseChangeMap_id}
  \uses{found:sym_power_base_change, found:modulecat_sym_power}
  Let \(R\) be a commutative ring, \(M\) an \(R\)-module and \(g:M\to M\) an
  \(R\)-linear map, viewed as an \(\mathrm{id}_R\)-semilinear map.  Then the
  base-change map \(\Sym^i g\) of \cref{found:sym_power_base_change} formed along
  the identity ring homomorphism \(\mathrm{id}_R\) coincides with the fixed-base
  symmetric power \(\Sym^i g\) of \cref{found:modulecat_sym_power}, both being
  \(R\)-linear endomorphisms of \(\Sym^i_R M\).
\end{lemma}

\begin{proof}\leanok
  \uses{found:sym_power_base_change, found:modulecat_sym_power}
  Both maps are \(R\)-linear endomorphisms of \(\Sym^i_R M\), so it suffices to
  compare them on the decomposable symmetric tensors \(m_1\cdots m_i\), which span.
  Base change along \(\mathrm{id}_R\) leaves the base ring unchanged, so its action
  \(m_1\cdots m_i\mapsto g(m_1)\cdots g(m_i)\) is exactly that of the fixed-base
  functor \cref{found:modulecat_sym_power}.
\end{proof}

\begin{lemma}[Base change respects composition of ring maps --- symmetric]\leanok
  \label{found:sym_base_change_map_comp}
  \lean{ModuleCat.symmetricPower.baseChangeMap_comp}
  \uses{found:sym_power_base_change}
  Let \(\varphi:R\to S\) and \(\psi:S\to T\) be homomorphisms of commutative rings,
  let \(g:M\to N\) be \(\varphi\)-semilinear and \(h:N\to P\) be
  \(\psi\)-semilinear.  Then the base-change map of the composite \(h\circ g\)
  (which is \((\psi\circ\varphi)\)-semilinear), formed along \(\psi\circ\varphi\),
  factors as the composite of the base-change maps of \(g\) and \(h\):
  \[
    \Sym^i(h\circ g)=\Sym^i h\circ\Sym^i g
    :\Sym^i_R M\to\Sym^i_T P .
  \]
\end{lemma}

\begin{proof}\leanok
  \uses{found:sym_power_base_change}
  Both sides are additive maps \(\Sym^i_R M\to\Sym^i_T P\) that agree on the
  spanning decomposable symmetric tensors: each sends \(m_1\cdots m_i\) to
  \(h(g(m_1))\cdots h(g(m_i))\) read over \(T\).  Hence they coincide.
\end{proof}

\begin{lemma}\leanok
[Generator action of the symmetric base-change map]
  \label{found:sym_base_change_map_tprod}
  \lean{ModuleCat.symmetricPower.baseChangeMap_tprod}
  \uses{found:sym_power_base_change, found:sympow_base_change_linear}
  The base-change map \cref{found:sym_power_base_change} along \(\varphi:R\to S\)
  sends the decomposable symmetric tensor \(m_1\cdots m_i\) over \(R\) to
  \(g(m_1)\cdots g(m_i)\) over \(S\).
\end{lemma}

\begin{proof}
  \uses{found:sym_power_base_change, found:sympow_base_change_linear}
  \leanok
  Unfolding \cref{found:sym_power_base_change} as the fixed-base symmetric power of
  \(g\) over \(R\) followed by the comparison map
  \cref{found:sympow_base_change_linear}, the tensor \(m_1\cdots m_i\) maps first
  to \(g(m_1)\cdots g(m_i)\) over \(R\) and then to the same tensor read over
  \(S\).
\end{proof}

\begin{lemma}\leanok
[Generator action of the symmetric-power comparison]
  \label{found:sympow_base_change_tprod}
  \lean{SymmetricPower.baseChange_tprod}
  \uses{found:sympow_base_change_linear, found:sympow_base_change_aux, found:sympow_base_change_aux_tprod}
  The comparison map \cref{found:sympow_base_change_linear} sends the decomposable
  symmetric tensor \(m_1\cdots m_i\) formed over \(A\) to the same tensor
  \(m_1\cdots m_i\) formed over \(B\).
\end{lemma}

\begin{proof}
  \uses{found:sympow_base_change_linear, found:sympow_base_change_aux, found:sympow_base_change_aux_tprod}
  \leanok
  By construction \cref{found:sympow_base_change_linear} descends the factorwise
  lift \cref{found:sympow_base_change_aux}; on the decomposable tensor
  \(m_1\cdots m_i\) the latter returns \(m_1\cdots m_i\) over \(B\)
  (\cref{found:sympow_base_change_aux_tprod}), which is therefore the value of the
  comparison map.
\end{proof}

\begin{lemma}\leanok
[Factorwise lift underlying the symmetric base-change comparison]
  \label{found:sympow_base_change_aux}
  \lean{SymmetricPower.baseChangeAux}
  \uses{mathlib:symmetricpower_type}
  The \(A\)-linear map \(\bigotimes^i_A N\to\Sym^i_B N\) from the \(i\)-fold
  \(A\)-tensor power obtained as the universal lift of the
  decomposable-symmetric-tensor family
  \((n_1,\dots,n_i)\mapsto n_1\cdots n_i\in\Sym^i_B N\), the latter's
  \(B\)-multilinearity restricted to \(A\).  It is the factorwise lift that
  descends to the comparison map \cref{found:sympow_base_change_linear}.
\end{lemma}

\begin{proof}
  \uses{mathlib:symmetricpower_type}
  \leanok
  The decomposable symmetric tensors of \(\Sym^i_B N\)
  (\cref{mathlib:symmetricpower_type}) assemble into a \(B\)-multilinear, hence
  \(A\)-multilinear, family; its universal lift through the \(A\)-tensor power is
  the stated \(A\)-linear map.
\end{proof}

\begin{lemma}\leanok
[The factorwise lift on a decomposable tensor]
  \label{found:sympow_base_change_aux_tprod}
  \lean{SymmetricPower.baseChangeAux_apply_tprod}
  \uses{mathlib:symmetricpower_type, found:sympow_base_change_aux}
  The factorwise lift \cref{found:sympow_base_change_aux} sends a decomposable
  tensor \(m_1\otimes\cdots\otimes m_i\) of the \(A\)-tensor power to the
  decomposable symmetric tensor \(m_1\cdots m_i\in\Sym^i_B N\).
\end{lemma}

\begin{proof}
  \uses{found:sympow_base_change_aux}
  \leanok
  This is the value of the universal lift on a decomposable tensor, namely the
  defining equation of \cref{found:sympow_base_change_aux}.
\end{proof}

\begin{lemma}\leanok
[The factorwise lift respects symmetrisation]
  \label{found:sympow_base_change_aux_respects}
  \lean{SymmetricPower.baseChangeAux_respects}
  \uses{mathlib:symmetricpower_type, found:sympow_base_change_aux, found:sympow_base_change_aux_tprod}
  The factorwise lift \cref{found:sympow_base_change_aux} is invariant under
  permutations of the tensor factors: for a permutation \(\sigma\), the
  decomposable tensors \(m_1\otimes\cdots\otimes m_i\) and
  \(m_{\sigma(1)}\otimes\cdots\otimes m_{\sigma(i)}\) have the same image.  Hence
  it descends through the \(A\)-symmetrisation quotient defining \(\Sym^i_A N\).
\end{lemma}

\begin{proof}
  \uses{found:sympow_base_change_aux, found:sympow_base_change_aux_tprod}
  \leanok
  By \cref{found:sympow_base_change_aux_tprod} both tensors map to a decomposable
  symmetric tensor over \(B\) that differ only by the permutation \(\sigma\) of
  their factors; since decomposable symmetric tensors over \(B\) are themselves
  permutation-invariant (\cref{mathlib:symmetricpower_type}), the two images
  coincide, so the lift is constant on \(A\)-symmetrisation classes.
\end{proof}

\begin{lemma}[Symmetric power of a presheaf of modules]\leanok
  \label{found:presheafmod_sym_power}
  \lean{AlgebraicGeometry.PresheafOfModules.symPower}
  \uses{found:modulecat_sym_power, found:ringcatsheaf_commring, found:sym_power_base_change, found:sym_base_change_map_id, found:sym_base_change_map_comp, found:sym_base_change_map_cat, found:presheafmod_sympowmap, found:presheafmod_sympowmap_id, found:presheafmod_sympowmap_comp}
  For \(i\in\mathbb N\), the \(i\)-th symmetric power functor
  \(\Sym^i:\mathrm{PresheafOfModules}\,R\to\mathrm{PresheafOfModules}\,R\), given
  over each open \(U\) by the module-level symmetric power
  \cref{found:modulecat_sym_power} of \(\mathcal F(U)\) and functorial in the
  restriction maps.
\end{lemma}

\begin{proof}\leanok
  \uses{found:modulecat_sym_power, found:ringcatsheaf_commring, found:sym_power_base_change, found:sym_base_change_map_id, found:sym_base_change_map_comp, found:sym_base_change_map_cat, found:presheafmod_sympowmap, found:presheafmod_sympowmap_id, found:presheafmod_sympowmap_comp}
  Two obstructions must be cleared, exactly as for the exterior power
  \cref{found:presheafmod_exterior_power}.  First,
  \(\mathrm{PresheafOfModules}\,R\) is \(\mathrm{RingCat}\)-valued, so the sections
  ring \(R(U)\) carries no commutative ring structure directly, and the
  module-level symmetric-power functor \cref{found:modulecat_sym_power}, which
  requires a commutative base, does not apply directly;
  \cref{found:ringcatsheaf_commring} restores the \(\mathrm{CommRing}\) structure
  on each \(R(U)\), so the openwise symmetric power \(\Sym^i_{R(U)}\mathcal F(U)\)
  is defined.  Second, the restriction map of \(\mathcal F\) along an inclusion
  \(V\subseteq U\) is not \(R(U)\)-linear but \emph{semilinear} over the ring map
  \(R(U)\to R(V)\), so the induced map on symmetric powers must run
  \(\Sym^i_{R(U)}\mathcal F(U)\to\Sym^i_{R(V)}\mathcal F(V)\) over
  \(R(U)\to R(V)\); this is exactly the base-change functoriality of
  \cref{found:sym_power_base_change}, not plain functoriality over a fixed base.
  Applying that lemma to each restriction map produces semilinear transition maps
  on symmetric powers, which inherit the presheaf identity and composition laws
  from those of \(\mathcal F\) --- via the identity-ring-map and composition
  naturality of base change, \cref{found:sym_base_change_map_id} and
  \cref{found:sym_base_change_map_comp} --- and assemble into a presheaf of modules
  \(\Sym^i\mathcal F\), functorially in \(\mathcal F\).
\end{proof}

\begin{definition}[Symmetric power of an \(\cO_X\)-module]\leanok
  \label{found:oxmod_sym_power}
  \lean{AlgebraicGeometry.Scheme.Modules.symPower}
  \uses{found:presheafmod_sym_power, mathlib:sheafification_modules_exact, mathlib:oxmod_category}
  For \(i\in\mathbb N\), the symmetric power functor
  \(\Sym^i:X.\mathrm{Modules}\to X.\mathrm{Modules}\), obtained by applying the
  presheaf-level symmetric power \cref{found:presheafmod_sym_power} to the
  underlying presheaf of an \(\cO_X\)-module and then sheafifying
  (\cref{mathlib:sheafification_modules_exact}).  Kemeny's \(\Sym^i\cS\) and
  \(\Sym^{k+1}\cG\) are instances.
\end{definition}

\subsection{Local freeness}

\begin{definition}[Locally free \(\cO_X\)-module]
  \leanok
  \label{found:locally_free}
  \lean{AlgebraicGeometry.Scheme.Modules.IsLocallyFree}
  \uses{mathlib:oxmod_category}
  The predicate \(\mathrm{IsLocallyFree}\) on \(X.\mathrm{Modules}\): an
  \(\cO_X\)-module \(\mathcal F\) is locally free (of rank \(r\)) when \(X\) is
  covered by open subsets on which \(\mathcal F\) is isomorphic to
  \(\cO_X^{\oplus r}\).  The rank is recorded as a finite index type \(I\) (of
  cardinality the rank), so this is local freeness of \emph{constant finite
  rank}.  Kemeny's ``locally free of rank \(k\)'' assertions for the cokernel
  bundle \(\mathcal W\) and the universal secant bundles are instances of this
  predicate.
\end{definition}

\subsection{Derived foundational anchors}

\begin{definition}[Higher direct image \(R^if_*\)]
  \leanok
  \label{def:peer_higher_direct_image}
  \lean{AlgebraicGeometry.higherDirectImage}
  \uses{mathlib:oxmod_pushforward, mathlib:right_derived, mathlib:short_exact, found:has_injective_resolutions}
  For a morphism of schemes \(f:X\to Y\) and \(i\in\mathbb N\), the higher direct
  image is the \(i\)-th right-derived functor of the additive pushforward:
  \[
    (R^if_*)\mathcal F=\bigl((\text{pushforward }f).\mathrm{rightDerived}\,i\bigr)\mathcal F .
  \]
  Right-deriving in the module category requires injective resolutions in
  \(X.\mathrm{Modules}\) (\cref{found:has_injective_resolutions}).  A short exact
  sequence of \(\cO_X\)-modules induces the long exact sequence of higher direct
  images.
\end{definition}

\begin{definition}[Sheaf cohomology \(H^i(X,\mathcal F)\)]
  \leanok
  \label{def:peer_scheme_hmodule}
  \lean{AlgebraicGeometry.Scheme.HModule}
  \uses{mathlib:right_derived, mathlib:short_exact, found:modules_to_absheaf, found:absheaf_global_sections, found:absheaf_global_sections_additive, found:absheaf_has_injective_resolutions}
  For a scheme \(X\) and \(i\in\mathbb N\), the sheaf cohomology of an
  \(\cO_X\)-module \(\mathcal F\) is computed in the category
  \(\mathrm{Sh}(X,\mathrm{Ab})\) of abelian-group sheaves: it is the \(i\)-th
  right-derived functor of global sections on \(\mathrm{Sh}(X,\mathrm{Ab})\),
  applied to the underlying abelian-group sheaf
  \(\mathcal F^{\mathrm{Ab}}\) (\cref{found:modules_to_absheaf}) of
  \(\mathcal F\):
  \[
    H^i(X,\mathcal F)=\bigl(\Gamma.\mathrm{rightDerived}\,i\bigr)\bigl(\mathcal F^{\mathrm{Ab}}\bigr).
  \]
  Injective resolutions in \(\mathrm{Sh}(X,\mathrm{Ab})\) are provided by
  \cref{found:absheaf_has_injective_resolutions}, requiring no hypothesis on
  \(X.\mathrm{Modules}\).  A short exact sequence of \(\cO_X\)-modules has exact
  underlying abelian-group sheaves, hence induces the long exact sequence of
  cohomology groups.
\end{definition}

\begin{definition}[Base-change comparison morphism]
  \leanok
  \label{def:peer_affine_pushforward_base_change}
  \lean{AlgebraicGeometry.pushforwardBaseChangeMap}
  \uses{mathlib:oxmod_pushforward}
  Given a commutative square of schemes
  \[
    \begin{array}{ccc}
      X' & \xrightarrow{\,g'\,} & X\\
      \downarrow{\scriptstyle f'} & & \downarrow{\scriptstyle f}\\
      Y' & \xrightarrow{\,g\,} & Y
    \end{array}
  \]
  the base-change comparison morphism is the natural map
  \[
    g^*f_*\mathcal F\to f'_*g'^*\mathcal F
  \]
  on degree-zero pushforwards, defined as the Beck--Chevalley mate of the
  canonical comparison of inverse-image functors under the two
  pullback--pushforward adjunctions.  Only commutativity of the square is needed
  to define the morphism; cartesianness of the square is the extra hypothesis that
  makes it an isomorphism.  Kemeny's pushforward identities for blow-ups, restrictions and products are
  expressed through this morphism.
\end{definition}

\begin{lemma}[Stalk-local criterion for module isomorphisms]
  \leanok
  \label{def:peer_module_iso_locality}
  \lean{AlgebraicGeometry.Modules.isIso_iff_isIso_stalkFunctor_map}
  \uses{mathlib:stalk_locality}
  A morphism of \(\cO_X\)-modules is an isomorphism if and only if the induced
  morphism on every stalk is an isomorphism.  The criterion reduces to the
  stalkwise criterion for morphisms of sheaves on the underlying topological
  space of \(X\).  A stronger criterion checks isomorphism on every open section
  of \(X\), which subsumes the stalkwise test.
\end{lemma}

\subsection{Coverage-debt helper functors}

\begin{lemma}\leanok
[Forgetful functor to abelian-group presheaves is additive]
  \label{found:modules_topresheaf_additive}
  \lean{AlgebraicGeometry.Scheme.Modules.toPresheaf_additive}
  \uses{mathlib:oxmod_category}
  For a scheme \(X\), the forgetful functor \(X.\mathrm{Modules}\to\) (presheaves
  of abelian groups on \(X\)) is additive.
\end{lemma}

\begin{proof}
  \uses{mathlib:oxmod_category}
  \leanok
  This forgetful functor is a composite of two additive functors (the forgetful
  from sheaves of modules to presheaves of modules, then the forgetful to
  presheaves of abelian groups), hence additive.
\end{proof}

\begin{lemma}\leanok
[Additive global-sections functor is additive]
  \label{found:global_sections_ab_additive}
  \lean{AlgebraicGeometry.Scheme.globalSectionsAb_additive}
  \uses{found:global_sections_ab, found:modules_topresheaf_additive}
  For a scheme \(X\), the global-sections functor
  \(\Gamma_{\mathrm{Ab}}\) of \cref{found:global_sections_ab} is additive.
\end{lemma}

\begin{proof}
  \uses{found:global_sections_ab, found:modules_topresheaf_additive}
  \leanok
  \(\Gamma_{\mathrm{Ab}}\) is the composite of the additive forgetful functor to
  abelian-group presheaves with the additive sections functor \(\Gamma(X,-)\); a
  composite of additive functors is additive.
\end{proof}

\begin{lemma}\leanok
[Categorical exterior base-change morphism]
  \label{found:exterior_base_change_map_cat}
  \lean{ModuleCat.exteriorPower.baseChangeMapCat}
  \uses{found:exterior_power_base_change, found:exterior_base_change_map_apply_iotamulti}
  For a homomorphism \(\varphi:R\to S\) of commutative rings, an \(R\)-module \(M\),
  an \(S\)-module \(N\) and a morphism \(g:M\to(\mathrm{restrictScalars}\,\varphi)\,N\)
  of \(\mathrm{Mod}_R\) (equivalently a \(\varphi\)-semilinear map \(M\to N\)), the
  induced morphism of \(\mathrm{Mod}_R\)
  \[
    \textstyle\bigwedge^i_R M\to(\mathrm{restrictScalars}\,\varphi)\bigl(\bigwedge^i_S N\bigr),
  \]
  the \(\mathrm{ModuleCat}\)-morphism repackaging of the semilinear base-change map
  \cref{found:exterior_power_base_change} across the additive bijection
  \(\mathrm{semilinearMapAddEquiv}\,\varphi\) between \(\varphi\)-semilinear maps
  \(M\to N\) and \(R\)-linear maps \(M\to(\mathrm{restrictScalars}\,\varphi)N\).  This
  is the form consumed by the openwise exterior-power presheaf assembly
  \cref{found:presheafmod_exterior_power}.
\end{lemma}

\begin{proof}
  \uses{found:exterior_power_base_change, found:exterior_base_change_map_apply_iotamulti}
  \leanok
  Apply \(\mathrm{semilinearMapAddEquiv}\,\varphi\) to the semilinear base-change map
  \(\wedge^i g\) of \cref{found:exterior_power_base_change}; its value on generators
  is read off from \cref{found:exterior_base_change_map_apply_iotamulti}.
\end{proof}

\begin{lemma}\leanok
[Generator action of the categorical exterior base-change morphism]
  \label{found:exterior_base_change_map_cat_hom_iotamulti}
  \lean{ModuleCat.exteriorPower.baseChangeMapCat_hom_ιMulti}
  \uses{found:exterior_base_change_map_cat}
  The morphism \cref{found:exterior_base_change_map_cat} along \(\varphi:R\to S\)
  sends a decomposable wedge \(v_1\wedge\cdots\wedge v_i\) over \(R\) to
  \(g(v_1)\wedge\cdots\wedge g(v_i)\) over \(S\).
\end{lemma}

\begin{proof}
  \uses{found:exterior_base_change_map_cat}
  \leanok
  Unfolding the \(\mathrm{semilinearMapAddEquiv}\) repackaging reduces the claim to
  the generator action of the underlying semilinear base-change map
  \cref{found:exterior_base_change_map_cat}.
\end{proof}

\begin{lemma}\leanok
[Categorical symmetric base-change morphism]
  \label{found:sym_base_change_map_cat}
  \lean{ModuleCat.symmetricPower.baseChangeMapCat}
  \uses{found:sym_power_base_change, found:sym_base_change_map_tprod}
  The symmetric twin of \cref{found:exterior_base_change_map_cat}: for
  \(\varphi:R\to S\), an \(R\)-module \(M\), an \(S\)-module \(N\) and a morphism
  \(g:M\to(\mathrm{restrictScalars}\,\varphi)\,N\), the induced morphism of
  \(\mathrm{Mod}_R\)
  \[
    \Sym^i_R M\to(\mathrm{restrictScalars}\,\varphi)\bigl(\Sym^i_S N\bigr),
  \]
  the \(\mathrm{ModuleCat}\)-morphism repackaging of the semilinear symmetric
  base-change map \cref{found:sym_power_base_change} across
  \(\mathrm{semilinearMapAddEquiv}\,\varphi\).  It is the form consumed by the
  openwise symmetric-power presheaf assembly \cref{found:presheafmod_sym_power}.
\end{lemma}

\begin{proof}
  \uses{found:sym_power_base_change, found:sym_base_change_map_tprod}
  \leanok
  Apply \(\mathrm{semilinearMapAddEquiv}\,\varphi\) to the semilinear base-change map
  \(\Sym^i g\) of \cref{found:sym_power_base_change}; its value on the spanning
  decomposable symmetric tensors is \cref{found:sym_base_change_map_tprod}.
\end{proof}

\begin{lemma}\leanok
[Generator action of the categorical symmetric base-change morphism]
  \label{found:sym_base_change_map_cat_tprod}
  \lean{ModuleCat.symmetricPower.baseChangeMapCat_hom_tprod}
  \uses{found:sym_base_change_map_cat}
  The morphism \cref{found:sym_base_change_map_cat} along \(\varphi:R\to S\) sends a
  decomposable symmetric tensor \(m_1\cdots m_i\) over \(R\) to
  \(g(m_1)\cdots g(m_i)\) over \(S\).
\end{lemma}

\begin{proof}
  \uses{found:sym_base_change_map_cat}
  \leanok
  Unfolding the \(\mathrm{semilinearMapAddEquiv}\) repackaging reduces the claim to
  the generator action of the underlying semilinear symmetric base-change map
  \cref{found:sym_base_change_map_cat}.
\end{proof}

\begin{lemma}\leanok
[Openwise exterior transition component]
  \label{found:presheafmod_extpowmap}
  \lean{AlgebraicGeometry.PresheafOfModules.extPowMap}
  \uses{found:exterior_base_change_map_cat}
  For a presheaf of modules \(\mathcal F\) over the structure ring presheaf of a
  scheme \(X\) and a morphism \(f\) of opens, the transition map on \(i\)-th exterior
  powers
  \(\bigwedge^i_{R(U)}\mathcal F(U)\to(\mathrm{restrictScalars})\bigwedge^i_{R(V)}\mathcal F(V)\),
  obtained by feeding the restriction ring map \(R(U)\to R(V)\) and the restriction
  \(\mathcal F.\mathrm{map}\,f\) to the categorical base-change morphism
  \cref{found:exterior_base_change_map_cat}.
\end{lemma}

\begin{proof}
  \uses{found:exterior_base_change_map_cat}
  \leanok
  Direct definition: \cref{found:exterior_base_change_map_cat} of the restriction ring
  map applied to \(\mathcal F.\mathrm{map}\,f\).
\end{proof}

\begin{lemma}\leanok
[Identity coherence of the exterior transition component]
  \label{found:presheafmod_extpowmap_id}
  \lean{AlgebraicGeometry.PresheafOfModules.extPowMap_id}
  \uses{found:presheafmod_extpowmap, found:exterior_base_change_map_cat_hom_iotamulti}
  The transition component \cref{found:presheafmod_extpowmap} along the identity map
  of an open is the identity, modulo the canonical restriction-of-scalars coherence
  isomorphism.
\end{lemma}

\begin{proof}
  \uses{found:presheafmod_extpowmap, found:exterior_base_change_map_cat_hom_iotamulti}
  \leanok
  Both sides agree on the spanning decomposable wedges by the generator action
  \cref{found:exterior_base_change_map_cat_hom_iotamulti} together with
  \(\mathcal F.\mathrm{map}\,\mathrm{id}=\mathrm{id}\).
\end{proof}

\begin{lemma}\leanok
[Composition coherence of the exterior transition component]
  \label{found:presheafmod_extpowmap_comp}
  \lean{AlgebraicGeometry.PresheafOfModules.extPowMap_comp}
  \uses{found:presheafmod_extpowmap, found:exterior_base_change_map_cat_hom_iotamulti}
  The transition component \cref{found:presheafmod_extpowmap} of a composite of opens
  is the composite of the components, modulo the restriction-of-scalars coherence.
\end{lemma}

\begin{proof}
  \uses{found:presheafmod_extpowmap, found:exterior_base_change_map_cat_hom_iotamulti}
  \leanok
  Both sides agree on the spanning decomposable wedges by the generator action
  \cref{found:exterior_base_change_map_cat_hom_iotamulti} together with
  \(\mathcal F.\mathrm{map}(f\circ g)=\mathcal F.\mathrm{map}\,g\circ\mathcal F.\mathrm{map}\,f\).
\end{proof}

\begin{lemma}\leanok
[Openwise symmetric transition component]
  \label{found:presheafmod_sympowmap}
  \lean{AlgebraicGeometry.PresheafOfModules.symPowMap}
  \uses{found:sym_base_change_map_cat}
  The symmetric twin of \cref{found:presheafmod_extpowmap}: for a presheaf of modules
  \(\mathcal F\) and a morphism \(f\) of opens, the transition map on \(i\)-th
  symmetric powers
  \(\Sym^i_{R(U)}\mathcal F(U)\to(\mathrm{restrictScalars})\Sym^i_{R(V)}\mathcal F(V)\),
  obtained by feeding the restriction ring map and \(\mathcal F.\mathrm{map}\,f\) to
  \cref{found:sym_base_change_map_cat}.
\end{lemma}

\begin{proof}
  \uses{found:sym_base_change_map_cat}
  \leanok
  Direct definition: \cref{found:sym_base_change_map_cat} of the restriction ring map
  applied to \(\mathcal F.\mathrm{map}\,f\).
\end{proof}

\begin{lemma}\leanok
[Identity coherence of the symmetric transition component]
  \label{found:presheafmod_sympowmap_id}
  \lean{AlgebraicGeometry.PresheafOfModules.symPowMap_id}
  \uses{found:presheafmod_sympowmap, found:sym_base_change_map_cat_tprod}
  The transition component \cref{found:presheafmod_sympowmap} along the identity map
  of an open is the identity, modulo the restriction-of-scalars coherence.
\end{lemma}

\begin{proof}
  \uses{found:presheafmod_sympowmap, found:sym_base_change_map_cat_tprod}
  \leanok
  Both sides agree on the spanning decomposable symmetric tensors by
  \cref{found:sym_base_change_map_cat_tprod} together with
  \(\mathcal F.\mathrm{map}\,\mathrm{id}=\mathrm{id}\).
\end{proof}

\begin{lemma}\leanok
[Composition coherence of the symmetric transition component]
  \label{found:presheafmod_sympowmap_comp}
  \lean{AlgebraicGeometry.PresheafOfModules.symPowMap_comp}
  \uses{found:presheafmod_sympowmap, found:sym_base_change_map_cat_tprod}
  The transition component \cref{found:presheafmod_sympowmap} of a composite of opens
  is the composite of the components, modulo the restriction-of-scalars coherence.
\end{lemma}

\begin{proof}
  \uses{found:presheafmod_sympowmap, found:sym_base_change_map_cat_tprod}
  \leanok
  Both sides agree on the spanning decomposable symmetric tensors by
  \cref{found:sym_base_change_map_cat_tprod} together with
  \(\mathcal F.\mathrm{map}(f\circ g)=\mathcal F.\mathrm{map}\,g\circ\mathcal F.\mathrm{map}\,f\).
\end{proof}

\subsection{Evaluation map and kernel bundle}

This subsection isolates the general foundational substrate behind Kemeny's kernel
bundle \(M_L\): for an arbitrary scheme \(X\) over a base ring \(k\) and an arbitrary
\(\cO_X\)-module \(\mathcal F\), the trivial bundle on the global sections, the
evaluation morphism onto \(\mathcal F\), and its kernel.  Specialising \(X\) to a K3
surface and \(\mathcal F\) to a line bundle \(L\) recovers Kemeny's defining sequence;
the curve-specific instantiation is taken up in the chapter on universal secant bundles.

\begin{definition}[Free sheaf of modules]
  \label{mathlib:sheafmod_free}
  \lean{SheafOfModules.free}
  \mathlibok
  \textit{Provided by Mathlib.}
  For a sheaf of rings \(R\) on a site and an index type \(I\), Mathlib provides the
  free sheaf of \(R\)-modules \(R^{(I)}=\mathrm{free}\,I\) on \(I\).
\end{definition}

\begin{lemma}[Universal property of the free sheaf of modules]
  \label{mathlib:sheafmod_freehomequiv}
  \lean{SheafOfModules.freeHomEquiv}
  \mathlibok
  \textit{Provided by Mathlib.}
  For a sheaf of \(R\)-modules \(\mathcal M\) and an index type \(I\), Mathlib
  provides the natural bijection
  \(\mathrm{Hom}(R^{(I)},\mathcal M)\cong(I\to\mathcal M.\mathrm{sections})\)
  identifying morphisms out of the free sheaf with \(I\)-indexed families of global
  sections --- the free \(\dashv\) global-sections adjunction at the level of
  hom-sets.
\end{lemma}

\begin{definition}[Categorical kernel]
  \label{mathlib:categorical_kernel}
  \lean{CategoryTheory.Limits.kernel}
  \mathlibok
  \textit{Provided by Mathlib.}
  In a category with zero morphisms and the relevant equalizer, Mathlib provides the
  kernel \(\ker f\) of a morphism \(f\) together with the monomorphism
  \(\ker f\to\mathrm{dom}\,f\) and the kernel's universal property.  In an abelian
  category every morphism has a kernel.
\end{definition}

\begin{definition}\leanok
[Trivial bundle on a vector space]
  \label{found:trivial_bundle}
  \lean{AlgebraicGeometry.Scheme.Modules.trivialBundle}
  \uses{mathlib:sheafmod_free}
  For a scheme \(X\) over a base ring \(k\) and a \(k\)-module \(V\), the trivial
  \(\cO_X\)-module \(V\otimes_k\cO_X\): the free \(\cO_X\)-module \(\cO_X^{(I)}\)
  (\cref{mathlib:sheafmod_free}) on an index set \(I\) of a basis of \(V\).  This is
  a construction only; when \(V\) is finite-dimensional of dimension \(r\) it is
  locally free of constant rank \(r\), recorded separately in
  \cref{found:trivial_bundle_locally_free}.
\end{definition}

\begin{proof}
  \uses{mathlib:sheafmod_free}
  \leanok
  Take a basis \(I\) of \(V\) and let \(V\otimes_k\cO_X:=\cO_X^{(I)}=\mathrm{free}\,I\).
  Local freeness of finite rank is \cref{found:trivial_bundle_locally_free}.
\end{proof}

\begin{lemma}\leanok
[Restriction of a free module along an open immersion is free]
  \label{found:restrict_free_iso}
  \lean{AlgebraicGeometry.Scheme.Modules.restrictFunctorObjFreeIso}
  \uses{found:trivial_bundle, mathlib:sheafmod_free}
  For an open immersion \(j:U\hookrightarrow X\) and an index type \(I\), the
  restriction \((\mathrm{restrictFunctor}\,j).\mathrm{obj}(\mathrm{free}\,I)\) of the
  free \(\cO_X\)-module on \(I\) is isomorphic to the free \(\cO_U\)-module
  \(\mathrm{free}\,I\).
\end{lemma}

\begin{proof}
  \uses{found:trivial_bundle, mathlib:sheafmod_free}
  \leanok
  The restriction functor along an open immersion is isomorphic to the inverse-image
  (pullback) functor \(\mathrm{restrictFunctor}\,j\cong\mathrm{pullback}\,j\), whose
  underlying site functor is the preimage \(\mathrm{Opens.map}\,j\).  For an open
  immersion this preimage is the right adjoint of the direct-image functor on opens,
  hence a final functor; finality is exactly the hypothesis under which pullback sends
  the free module \(\mathrm{free}\,I\) to \(\mathrm{free}\,I\).  Composing the two
  isomorphisms gives the claim.
\end{proof}

\begin{lemma}[Triviality descends along a scheme isomorphism]
  \label{found:free_iso_of_restrict_iso}
  \lean{AlgebraicGeometry.Scheme.Modules.freeIso_of_restrictFunctor_hom_freeIso}
  \uses{found:restrict_free_iso, mathlib:sheafmod_free}
  Let \(e:S\cong T\) be an isomorphism of schemes, \(A\) an \(\cO_T\)-module, and
  \(I\) an index type.  If the restriction of \(A\) along \(e\),
  \((\mathrm{restrictFunctor}\,e).\mathrm{obj}\,A\), is isomorphic to the free
  \(\cO_S\)-module \(\mathrm{free}\,I\) (\cref{mathlib:sheafmod_free}), then \(A\) is
  itself isomorphic to the free \(\cO_T\)-module \(\mathrm{free}\,I\).
\end{lemma}

\begin{proof}
  \uses{found:restrict_free_iso, mathlib:sheafmod_free}
  Restriction along the isomorphism \(e\) is invertible: composing the restriction
  functors along \(e\) and \(e^{-1}\) is naturally isomorphic to the identity (using
  the composition and identity coherences of \(\mathrm{restrictFunctor}\) together
  with \(e^{-1}\circ e=\mathrm{id}\)).  Applying \(\mathrm{restrictFunctor}\,e^{-1}\)
  to the hypothesis and using \(\mathrm{restrictFunctor}\,e^{-1}(\mathrm{free}\,I)
  \cong\mathrm{free}\,I\) (\cref{found:restrict_free_iso}, since \(e^{-1}\) is an open
  immersion) yields \(A\cong\mathrm{free}\,I\).
\end{proof}

\begin{lemma}\leanok
[Local freeness is invariant under isomorphism]
  \label{found:locally_free_of_iso}
  \lean{AlgebraicGeometry.Scheme.Modules.isLocallyFree_of_iso}
  \uses{found:locally_free}
  If \(\mathcal F\cong\mathcal G\) in \(X.\mathrm{Modules}\) and \(\mathcal F\) is
  locally free (\cref{found:locally_free}), then \(\mathcal G\) is locally free, with the
  same trivialising cover and the same finite rank.
\end{lemma}

\begin{proof}
  \uses{found:locally_free}
  \leanok
  Given a trivialising open \(U\) with \((\mathrm{restrictFunctor}\,U).\mathrm{obj}\,
  \mathcal F\cong\mathrm{free}\,I\), transport it along the restriction
  \((\mathrm{restrictFunctor}\,U).\mathrm{mapIso}\) of the isomorphism
  \(\mathcal F\cong\mathcal G\) to obtain a trivialisation of \(\mathcal G\) on the same
  open and the same finite index type \(I\); the trivialising cover is unchanged.
\end{proof}

\begin{lemma}\leanok
[Trivial bundle of finite rank is locally free]
  \label{found:trivial_bundle_locally_free}
  \lean{AlgebraicGeometry.Scheme.Modules.trivialBundle_isLocallyFree}
  \uses{found:trivial_bundle, found:locally_free, found:restrict_free_iso, mathlib:sheafmod_free}
  For a scheme \(X\) and a finite index type \(I\), the trivial bundle
  \(\mathrm{trivialBundle}\,I=\cO_X^{(I)}\) (\cref{found:trivial_bundle}) is locally
  free of constant rank \(|I|\) (\cref{found:locally_free}).  In particular, for a
  finite-dimensional \(k\)-vector space \(V\) of dimension \(r\) the trivial bundle
  \(V\otimes_k\cO_X\) is locally free of rank \(r\).
\end{lemma}

\begin{proof}
  \uses{found:trivial_bundle, found:locally_free, found:restrict_free_iso, mathlib:sheafmod_free}
  \leanok
  The free \(\cO_X\)-module \(\cO_X^{(I)}=\mathrm{free}\,I\) is already globally a
  direct sum of \(|I|\) copies of \(\cO_X\).  Taking the single trivialising open
  cover by \(U=X\) itself, the restriction of \(\cO_X^{(I)}\) to \(U\) is isomorphic
  to the free \(\cO_U\)-module \(\mathrm{free}\,I\), so every point of \(X\) lies in
  an open on which the bundle is free on the fixed finite index set \(I\); this is
  exactly local freeness of constant rank \(|I|\) in the sense of
  \cref{found:locally_free}.  For \(V\) finite-dimensional one applies this with
  \(I\) a basis of \(V\), of cardinality \(\dim_k V\).
\end{proof}

\begin{definition}\leanok
[Evaluation morphism]
  \label{found:eval_morphism}
  \lean{AlgebraicGeometry.Scheme.Modules.evalMap}
  \uses{found:trivial_bundle, found:absheaf_global_sections, mathlib:sheafmod_freehomequiv}
  For an \(\cO_X\)-module \(\mathcal F\), the evaluation morphism
  \[
    \mathrm{ev}_{\mathcal F}:H^0(X,\mathcal F)\otimes_k\cO_X\to\mathcal F
  \]
  out of the trivial bundle (\cref{found:trivial_bundle}) on the global sections
  \(H^0(X,\mathcal F)\) (\cref{found:absheaf_global_sections}): the
  \(\cO_X\)-linear extension of \(1\mapsto s\) on each basis section
  \(s\in H^0(X,\mathcal F)\).  Equivalently it is the counit of the adjunction between
  the free-\(\cO_X\)-module functor and the global-sections functor, corresponding
  under the free hom-bijection \cref{mathlib:sheafmod_freehomequiv} to the identity
  family of global sections.
\end{definition}

\begin{proof}
  \uses{found:trivial_bundle, found:absheaf_global_sections, mathlib:sheafmod_freehomequiv}
  \leanok
  A global section \(s\in H^0(X,\mathcal F)\) is a morphism \(\cO_X\to\mathcal F\)
  with \(1\mapsto s\); by the free hom-bijection \cref{mathlib:sheafmod_freehomequiv}
  the family \((s)_{s}\) of all global sections assembles into a single morphism out
  of the free bundle \(H^0(X,\mathcal F)\otimes_k\cO_X\), natural in \(\mathcal F\)
  (the counit).
\end{proof}

% SOURCE: Kemeny, Universal Secant Bundles and Syzygies of Canonical Curves, \S1 (sec1), opening (read from references/MR4213770-secant-syzygies.tex, line 255)
% SOURCE QUOTE: "Let $X$ be a K3 surface of Picard rank one and even genus $g=2k$. Let $M_L$ denote the bundle defined by the sequence $0 \to M_L \to \mathrm{H}^0(L)\otimes \mathcal{O}_X \to L \to 0.$
% By the Kernel Bundle description of Koszul cohomology, see \cite{lazarsfeld-VBT} or \cite[\S 3]{ein-lazarsfeld-asymptotic}, to prove Voisin's Theorem it suffices to show $$\mathrm{H}^1(X, \bigwedge^{k+1} M_L )=0.$$"
\begin{definition}\leanok
[Evaluation kernel bundle]
  \label{found:eval_kernel_bundle}
  \lean{AlgebraicGeometry.Scheme.Modules.evalKernel}
  \uses{found:eval_morphism, mathlib:categorical_kernel, mathlib:oxmod_category}
  \textit{Source: Kemeny, \S1.}
  For an \(\cO_X\)-module \(\mathcal F\), the evaluation kernel
  \(M_{\mathcal F}:=\ker(\mathrm{ev}_{\mathcal F})\), the kernel
  (\cref{mathlib:categorical_kernel}) of the evaluation morphism
  \cref{found:eval_morphism} taken in the abelian category \(X.\mathrm{Modules}\)
  (\cref{mathlib:oxmod_category}).  It sits in the left-exact sequence
  \[
    0\to M_{\mathcal F}\to H^0(X,\mathcal F)\otimes_k\cO_X
       \xrightarrow{\mathrm{ev}_{\mathcal F}}\mathcal F .
  \]
  Specialising \(X\) to a K3 surface and \(\mathcal F\) to a line bundle \(L\) gives
  Kemeny's kernel bundle \(M_L\), defined by
  \(0\to M_L\to H^0(L)\otimes\cO_X\to L\to 0\).
\end{definition}

\begin{proof}
  \uses{found:eval_morphism, mathlib:categorical_kernel, mathlib:oxmod_category}
  \leanok
  The category \(X.\mathrm{Modules}\) is abelian (\cref{mathlib:oxmod_category}), so
  the morphism \(\mathrm{ev}_{\mathcal F}\) of \cref{found:eval_morphism} has a kernel
  (\cref{mathlib:categorical_kernel}); \(M_{\mathcal F}\) is by definition that kernel,
  with its monomorphism into \(H^0(X,\mathcal F)\otimes_k\cO_X\).
\end{proof}

\begin{definition}\leanok
[Evaluation short complex]
  \label{found:eval_short_complex}
  \lean{AlgebraicGeometry.Scheme.Modules.evalShortComplex}
  \uses{found:eval_morphism, found:eval_kernel_bundle, mathlib:categorical_kernel}
  For an \(\cO_X\)-module \(\mathcal F\), the short complex
  \[
    M_{\mathcal F}\hookrightarrow H^0(X,\mathcal F)\otimes_k\cO_X
       \xrightarrow{\mathrm{ev}_{\mathcal F}}\mathcal F ,
  \]
  whose first map is the kernel inclusion of the evaluation morphism
  (\cref{found:eval_morphism}) and whose second map is \(\mathrm{ev}_{\mathcal F}\); its
  left object is (definitionally) the evaluation kernel bundle \(M_{\mathcal F}\)
  (\cref{found:eval_kernel_bundle}).
\end{definition}

\begin{proof}
  \uses{found:eval_morphism, found:eval_kernel_bundle, mathlib:categorical_kernel}
  \leanok
  The short complex is assembled from the kernel inclusion
  \(\ker(\mathrm{ev}_{\mathcal F})\hookrightarrow H^0\otimes\cO_X\)
  (\cref{mathlib:categorical_kernel}) and \(\mathrm{ev}_{\mathcal F}\) itself; the
  composite is zero by the kernel condition.  Its left object is by definition
  \(M_{\mathcal F}=\ker(\mathrm{ev}_{\mathcal F})\) (\cref{found:eval_kernel_bundle}).
\end{proof}

\begin{lemma}\leanok
[Epimorphic evaluation gives a short exact evaluation sequence]
  \label{found:eval_shortexact_of_globally_generated}
  \lean{AlgebraicGeometry.Scheme.Modules.evalShortComplex_shortExact_of_globallyGenerated}
  \uses{found:eval_morphism, found:eval_kernel_bundle, found:eval_short_complex, mathlib:short_exact}
  If the evaluation morphism \(\mathrm{ev}_{\mathcal F}\) (\cref{found:eval_morphism}) is
  an epimorphism, then the evaluation short complex (\cref{found:eval_short_complex})
  \[
    0\to M_{\mathcal F}\to H^0(X,\mathcal F)\otimes_k\cO_X
       \xrightarrow{\mathrm{ev}_{\mathcal F}}\mathcal F\to 0
  \]
  is short exact (\cref{mathlib:short_exact}).  This is the formalized core of the
  paper's evaluation sequence: a globally generated \(\mathcal F\) has surjective
  evaluation (the global sections generate every stalk), so the epimorphism hypothesis
  holds.  The implication ``\(\mathcal F\) globally generated \(\Rightarrow\)
  \(\mathrm{ev}_{\mathcal F}\) epi'' is itself not yet formalized --- no
  \(\mathrm{GloballyGenerated}\) predicate on \(X.\mathrm{Modules}\) is available ---
  and is left for future work.
\end{lemma}

\begin{proof}
  \uses{found:eval_kernel_bundle, found:eval_short_complex, mathlib:short_exact}
  \leanok
  By construction \(M_{\mathcal F}=\ker(\mathrm{ev}_{\mathcal F})\)
  (\cref{found:eval_kernel_bundle}): the kernel inclusion is a monomorphism and the
  evaluation short complex (\cref{found:eval_short_complex}) is exact at the middle
  term.  The hypothesis that \(\mathrm{ev}_{\mathcal F}\) is an epimorphism upgrades this
  left-exact kernel sequence to a short exact sequence (\cref{mathlib:short_exact}).
\end{proof}

% --- Mathlib dependency anchors for the kernel-bundle locally-free decomposition
%     (\cref{found:isLocallyFree_kernel_of_shortExact}). ---

\begin{definition}[Tilde: quasicoherent sheaf of a module on \(\operatorname{Spec}R\)]
  \label{mathlib:tilde}
  \lean{AlgebraicGeometry.tilde}
  \mathlibok
  \textit{Provided by Mathlib.}
  For a commutative ring \(R\) and an \(R\)-module \(M\), Mathlib provides the
  associated quasicoherent sheaf of \(\cO_{\operatorname{Spec}R}\)-modules
  \(\widetilde M=\mathrm{tilde}\,M\) on \(\operatorname{Spec}R\), the
  module-to-sheaf construction \(M\mapsto M^{\sim}\).  Mathlib provides that it is
  additive and fully faithful (\cref{mathlib:tilde_fullyFaithful}) and is the left
  adjoint of the global-sections functor; as a left adjoint it preserves coproducts
  and cokernels, and it identifies \(R\)-modules with quasicoherent
  \(\cO_{\operatorname{Spec}R}\)-modules.  Mathlib does \emph{not} provide that
  \(\mathrm{tilde}\) preserves arbitrary kernels (left-exactness); the affine bridge
  (\cref{lem:kernel_iso_tilde_affine}) avoids needing it by exploiting that the
  relevant sequence \emph{splits}, on which an additive functor is automatically
  exact (\cref{mathlib:splitting_map}).
\end{definition}

\begin{lemma}[Tilde of a finite free module is the free sheaf]
  \label{mathlib:tildeFinsupp}
  \lean{AlgebraicGeometry.tildeFinsupp}
  \mathlibok
  \textit{Provided by Mathlib.}
  For a commutative ring \(R\) and an index type \(\iota\), Mathlib provides a
  natural isomorphism \(\widetilde{(\iota\to_0 R)}\cong\mathrm{free}\,\iota\)
  (\cref{mathlib:tilde}, \cref{mathlib:sheafmod_free}) between the tilde of the
  finitely supported free \(R\)-module on \(\iota\) and the free sheaf of
  \(\cO_{\operatorname{Spec}R}\)-modules on \(\iota\).
\end{lemma}

\begin{lemma}[Tilde is fully faithful]
  \label{mathlib:tilde_fullyFaithful}
  \lean{AlgebraicGeometry.tilde.fullyFaithfulFunctor}
  \mathlibok
  \textit{Provided by Mathlib.}
  Mathlib provides that the tilde functor (\cref{mathlib:tilde}) from
  \(R\)-modules to \(\cO_{\operatorname{Spec}R}\)-modules is fully faithful: the
  induced map on hom-sets
  \(\mathrm{Hom}_R(M,N)\to\mathrm{Hom}(\widetilde M,\widetilde N)\) is a bijection.
\end{lemma}

\begin{lemma}[A short exact sequence with projective quotient splits]
  \label{mathlib:splittingOfProjective}
  \lean{CategoryTheory.ShortComplex.ShortExact.splittingOfProjective}
  \mathlibok
  \textit{Provided by Mathlib.}
  In a balanced preadditive category, a short exact sequence
  \(0\to X_1\to X_2\to X_3\to 0\) (\cref{mathlib:short_exact}) whose right-hand
  term \(X_3\) is projective admits a splitting; in particular
  \(X_2\cong X_1\oplus X_3\).
\end{lemma}

\begin{lemma}[An additive functor carries a splitting to a splitting]
  \label{mathlib:splitting_map}
  \lean{CategoryTheory.ShortComplex.Splitting.map}
  \mathlibok
  \textit{Provided by Mathlib.}
  If \(s\) is a splitting of a short complex \(S=(X_1\to X_2\to X_3)\) and
  \(F\) is an additive functor, then \(F\) carries \(s\) to a splitting of the
  mapped short complex \(F(S)=(FX_1\to FX_2\to FX_3)\).  Consequently \(F(S)\) is
  short exact (\cref{mathlib:short_exact}) and the mapped first map \(FX_1\to FX_2\)
  is a kernel of \(FX_2\to FX_3\) — an additive functor is exact on split sequences,
  with no left-exactness hypothesis on \(F\).
\end{lemma}

\begin{definition}[Free locus of a module]
  \label{mathlib:freeLocus}
  \lean{Module.freeLocus}
  \mathlibok
  \textit{Provided by Mathlib.}
  For a commutative ring \(R\) and an \(R\)-module \(M\), Mathlib provides the free
  locus \(\mathrm{freeLocus}\,R\,M\subseteq\operatorname{Spec}R\), the set of primes
  \(\mathfrak p\) at which the localization \(M_{\mathfrak p}\) is free of finite
  rank.
\end{definition}

\begin{lemma}[The free locus is open for a finitely presented module]
  \label{mathlib:isOpen_freeLocus}
  \lean{Module.isOpen_freeLocus}
  \mathlibok
  \textit{Provided by Mathlib.}
  For a finitely presented \(R\)-module \(M\), the free locus
  \(\mathrm{freeLocus}\,R\,M\) (\cref{mathlib:freeLocus}) is open in
  \(\operatorname{Spec}R\).
\end{lemma}

\begin{lemma}[Free locus is everything iff projective]
  \label{mathlib:freeLocus_eq_univ_iff}
  \lean{Module.freeLocus_eq_univ_iff}
  \mathlibok
  \textit{Provided by Mathlib.}
  For a finitely presented \(R\)-module \(M\),
  \(\mathrm{freeLocus}\,R\,M=\operatorname{Spec}R\) if and only if \(M\) is
  projective (\cref{mathlib:freeLocus}).
\end{lemma}

\begin{lemma}[Basic open in the free locus iff localized module projective]
  \label{mathlib:basicOpen_subset_freeLocus_iff}
  \lean{Module.basicOpen_subset_freeLocus_iff}
  \mathlibok
  \textit{Provided by Mathlib.}
  For a finitely presented \(R\)-module \(M\) and \(f\in R\), the basic open
  \(D(f)\) is contained in \(\mathrm{freeLocus}\,R\,M\) (\cref{mathlib:freeLocus})
  if and only if the localized module \(M_f\) is projective over \(R_f\); for a
  finitely presented module such a localized projective module is free of finite
  rank.
\end{lemma}

% --- Affine reduction sub-lemmas for \cref{found:isLocallyFree_kernel_of_shortExact}.
%     The Kemeny/Hartshorne source quote lives on that lemma block below. ---

\begin{lemma}\leanok
[Local freeness is local on the base]
  \label{lem:isLocallyFree_local}
  \lean{AlgebraicGeometry.Scheme.Modules.isLocallyFree_of_isLocallyFree_cover}
  \uses{found:locally_free, found:free_iso_of_restrict_iso}
  Let \(\{U_\alpha\}\) be an open cover of \(X\) and \(\mathcal K\) an
  \(\cO_X\)-module.  If each restriction \(\mathcal K|_{U_\alpha}\) is locally free
  of constant rank \(r\) (\cref{found:locally_free}), then \(\mathcal K\) is locally
  free of rank \(r\) on \(X\).
\end{lemma}

\begin{proof}
  \uses{found:locally_free, found:free_iso_of_restrict_iso}
  \leanok
  Local freeness of rank \(r\) means \(X\) has an open cover trivialising
  \(\mathcal K\) to \(\mathrm{free}\,I\) with \(|I|=r\) (\cref{found:locally_free}).
  Fix \(x\in X\); it lies in some \(U_\alpha\), and the trivialising cover of
  \(\mathcal K|_{U_\alpha}\) gives an open \(V\subseteq U_\alpha\) with
  \(x\in V\) and \((\mathcal K|_{U_\alpha})|_V\cong\mathrm{free}\,I\).  Pushing \(V\)
  forward to its open image \(W=j(V)\subseteq X\) along the open immersion
  \(j:U_\alpha\hookrightarrow X\) gives a scheme isomorphism \(e:V\cong W\) compatible
  with the inclusions, so that \((\mathrm{restrictFunctor}\,e)\) carries
  \(\mathcal K|_W\) to \((\mathcal K|_{U_\alpha})|_V\cong\mathrm{free}\,I\).  By
  \cref{found:free_iso_of_restrict_iso} this descends to a trivialisation
  \(\mathcal K|_W\cong\mathrm{free}\,I\) on the \(X\)-open \(W\ni x\).  As \(x\) was
  arbitrary, these \(W\) form an open cover of \(X\) trivialising \(\mathcal K\) on
  the fixed \(I\).
\end{proof}

\begin{lemma}\leanok
[Kernel of a finite free presentation is free on a basic-open cover]
  \label{lem:module_kernel_free_basicOpen_cover}
  \lean{AlgebraicGeometry.Scheme.Modules.kernel_module_free_on_basicOpen_cover}
  \uses{mathlib:short_exact, mathlib:splittingOfProjective, mathlib:freeLocus, mathlib:isOpen_freeLocus, mathlib:freeLocus_eq_univ_iff, mathlib:basicOpen_subset_freeLocus_iff}
  Let \(0\to K\to R^m\to R^n\to 0\) be a short exact sequence
  (\cref{mathlib:short_exact}) of \(R\)-modules with \(R^m\) and \(R^n\) finite
  free.  Then \(K\) is finitely presented and projective, and there is a finite
  family \(f_1,\dots,f_s\in R\) with \(\bigcup_i D(f_i)=\operatorname{Spec}R\) such
  that each localized module \(K_{f_i}\) is free of finite rank over \(R_{f_i}\).
\end{lemma}

\begin{proof}
  \uses{mathlib:short_exact, mathlib:splittingOfProjective, mathlib:freeLocus, mathlib:isOpen_freeLocus, mathlib:freeLocus_eq_univ_iff, mathlib:basicOpen_subset_freeLocus_iff}
  \leanok
  \(R^n\) is free, hence projective, so the sequence splits
  (\cref{mathlib:splittingOfProjective}) and \(R^m\cong K\oplus R^n\); thus \(K\) is a
  finitely presented projective summand of a finite free module.  Projectivity gives
  \(\mathrm{freeLocus}\,R\,K=\operatorname{Spec}R\)
  (\cref{mathlib:freeLocus_eq_univ_iff}, \cref{mathlib:freeLocus}).  This locus is
  open (\cref{mathlib:isOpen_freeLocus}) and the basic opens form a basis, so every
  point lies in some \(D(f)\subseteq\mathrm{freeLocus}\,R\,K\), on which \(K_f\) is
  projective, hence free of finite rank, over \(R_f\)
  (\cref{mathlib:basicOpen_subset_freeLocus_iff}).  Quasi-compactness of
  \(\operatorname{Spec}R\) extracts a finite subcover \(D(f_1),\dots,D(f_s)\).
\end{proof}

\begin{lemma}[Affine bridge: the kernel of a free SES on \(\operatorname{Spec}R\) is a tilde]
  \label{lem:kernel_iso_tilde_affine}
  \lean{AlgebraicGeometry.Scheme.Modules.kernel_iso_tilde_of_free_ses_affine}
  \uses{mathlib:short_exact, mathlib:tilde, mathlib:tildeFinsupp, mathlib:tilde_fullyFaithful, mathlib:sheafmod_free, mathlib:splittingOfProjective, mathlib:splitting_map}
  Let \(0\to\mathcal K\to\mathcal G\to\mathcal H\to 0\) be a short exact sequence
  (\cref{mathlib:short_exact}) of \(\cO_{\operatorname{Spec}R}\)-modules with
  \(\mathcal G\cong\mathrm{free}\,m\) and \(\mathcal H\cong\mathrm{free}\,n\)
  (\cref{mathlib:sheafmod_free}).  Then there is a short exact sequence of
  \(R\)-modules \(0\to K\to R^m\to R^n\to 0\) (\cref{mathlib:short_exact}) together
  with an isomorphism \(\mathcal K\cong\widetilde K\) (\cref{mathlib:tilde}).
\end{lemma}

\begin{proof}
  \uses{mathlib:short_exact, mathlib:tilde, mathlib:tildeFinsupp, mathlib:tilde_fullyFaithful, mathlib:sheafmod_free, mathlib:splittingOfProjective, mathlib:splitting_map}
  By \cref{mathlib:tildeFinsupp}, \(\mathrm{free}\,m\cong\widetilde{R^m}\) and
  \(\mathrm{free}\,n\cong\widetilde{R^n}\), so transporting the surjection
  \(\mathcal G\to\mathcal H\) along these isomorphisms presents it as a map
  \(\widetilde{R^m}\to\widetilde{R^n}\).  Full faithfulness
  (\cref{mathlib:tilde_fullyFaithful}) writes this map as \(\widetilde\varphi\) for a
  unique \(R\)-linear \(\varphi:R^m\to R^n\).  Since \(\mathcal G\to\mathcal H\) is an
  epimorphism and \(\mathrm{tilde}\) is faithful, hence reflects epimorphisms,
  \(\varphi\) is surjective; put \(K=\ker\varphi\), giving the \(R\)-module short
  exact sequence \(0\to K\to R^m\to R^n\to 0\) (\cref{mathlib:short_exact}).

  Crucially we do \emph{not} need \(\mathrm{tilde}\) to preserve kernels in general.
  The module sequence \emph{splits}: \(R^n\) is free, hence projective, so
  \(0\to K\to R^m\to R^n\to 0\) admits a splitting \(s\)
  (\cref{mathlib:splittingOfProjective}).  Because \(\mathrm{tilde}\) is additive,
  it carries \(s\) to a splitting of \(0\to\widetilde K\to\widetilde{R^m}\to
  \widetilde{R^n}\to 0\) (\cref{mathlib:splitting_map}); in particular the mapped
  inclusion \(\widetilde K\to\widetilde{R^m}\) is a kernel of \(\widetilde\varphi\).
  But \(\mathcal K\to\mathcal G\) is a kernel of \(\mathcal G\to\mathcal H\) (it is
  the kernel term of the given short exact sequence), and under the trivialising
  isomorphisms \(\mathcal G\cong\widetilde{R^m}\), \(\mathcal H\cong\widetilde{R^n}\)
  the map \(\mathcal G\to\mathcal H\) is identified with \(\widetilde\varphi\).  Two
  kernels of the same morphism are canonically isomorphic, so
  \(\mathcal K\cong\widetilde K\) (\cref{mathlib:tilde}).
\end{proof}

\begin{lemma}[Kernel is locally free on an affine trivialising open]
  \label{lem:isLocallyFree_kernel_on_affine}
  \lean{AlgebraicGeometry.Scheme.Modules.isLocallyFree_kernel_on_affine}
  \uses{lem:kernel_iso_tilde_affine, lem:module_kernel_free_basicOpen_cover, mathlib:tilde, mathlib:tildeFinsupp, mathlib:sheafmod_free, found:locally_free, found:locally_free_of_iso}
  Let \(0\to\mathcal K\to\mathcal G\to\mathcal H\to 0\) be a short exact sequence
  of \(\cO_{\operatorname{Spec}R}\)-modules with \(\mathcal G\cong\mathrm{free}\,m\)
  and \(\mathcal H\cong\mathrm{free}\,n\).  Then \(\mathcal K\) is locally free of
  constant rank \(m-n\) on \(\operatorname{Spec}R\) (\cref{found:locally_free}).
\end{lemma}

\begin{proof}
  \uses{lem:kernel_iso_tilde_affine, lem:module_kernel_free_basicOpen_cover, mathlib:tilde, mathlib:tildeFinsupp, mathlib:sheafmod_free, found:locally_free, found:locally_free_of_iso}
  The affine bridge (\cref{lem:kernel_iso_tilde_affine}) gives
  \(0\to K\to R^m\to R^n\to 0\) with \(\mathcal K\cong\widetilde K\)
  (\cref{mathlib:tilde}).  By \cref{lem:module_kernel_free_basicOpen_cover},
  \(\operatorname{Spec}R\) is covered by basic opens \(D(f_i)\) with \(K_{f_i}\) free
  of rank \(m-n\).  On \(D(f_i)\cong\operatorname{Spec}R_{f_i}\), \(\widetilde K\)
  restricts to \(\widetilde{K_{f_i}}\cong\mathrm{free}\,(m-n)\)
  (\cref{mathlib:tildeFinsupp}, \cref{mathlib:sheafmod_free}); transporting along
  \(\mathcal K\cong\widetilde K\) (\cref{found:locally_free_of_iso}) trivialises
  \(\mathcal K\) on each \(D(f_i)\).  These cover \(\operatorname{Spec}R\), so
  \(\mathcal K\) is locally free of rank \(m-n\) (\cref{found:locally_free}).
\end{proof}

% SOURCE: Kemeny, Universal Secant Bundles and Syzygies of Canonical Curves, \S0 (Introduction, preliminaries) (read from references/MR4213770-secant-syzygies.tex, line 251)
% SOURCE QUOTE: "Assume we have an exact sequence $0 \to \mathcal{F} \to \mathcal{G} \to \mathcal{H} \to 0$ of coherent sheaves on a quasi-projective variety, with $\mathcal{G}$ locally free. Assume either $\mathcal{H}$ is locally free or $\mathcal{H} \simeq \mathcal{O}_D$ for a Cartier divisor $D$. Then $\mathcal{F}$ is locally free. This follows from \cite[III, Ex 6.5]{hartshorne}."
\begin{lemma}[Locally free kernel of a short exact sequence (Hartshorne III, Ex.~6.5)]
  \label{found:isLocallyFree_kernel_of_shortExact}
  \lean{AlgebraicGeometry.Scheme.Modules.isLocallyFree_kernel_of_shortExact}
  \uses{found:locally_free, mathlib:short_exact, lem:isLocallyFree_local, lem:isLocallyFree_kernel_on_affine}
  \textit{Source: Kemeny, \S0 (preliminaries on locally free sheaves), citing
  Hartshorne III, Ex.~6.5.}
  Let \(0\to\mathcal K\to\mathcal G\to\mathcal H\to 0\) be a short exact sequence
  (\cref{mathlib:short_exact}) of \(\cO_X\)-modules with \(\mathcal G\) and
  \(\mathcal H\) locally free (\cref{found:locally_free}).  Then \(\mathcal K\) is
  locally free.
\end{lemma}

\begin{proof}
  \uses{lem:isLocallyFree_local, lem:isLocallyFree_kernel_on_affine, found:locally_free, mathlib:short_exact}
  Local freeness is local on \(X\) (\cref{lem:isLocallyFree_local}).  Each point has
  an affine neighbourhood \(\operatorname{Spec}R\subseteq X\) trivialising both
  \(\mathcal G\) and \(\mathcal H\) (refine an affine cover against their
  trivialising covers, \cref{found:locally_free}).  There the sequence
  (\cref{mathlib:short_exact}) has free middle and right terms, so
  \cref{lem:isLocallyFree_kernel_on_affine} trivialises \(\mathcal K\).  These opens
  cover \(X\), so \(\mathcal K\) is locally free (\cref{lem:isLocallyFree_local}).
  This is Hartshorne III, Ex.~6.5.
\end{proof}

\begin{lemma}[Rank additivity for a short exact sequence of locally free sheaves]
  \label{found:rank_kernel_of_shortExact}
  \lean{AlgebraicGeometry.Scheme.Modules.rank_kernel_of_shortExact}
  \uses{found:locally_free, found:isLocallyFree_kernel_of_shortExact, mathlib:short_exact}
  Let \(0\to\mathcal K\to\mathcal G\to\mathcal H\to 0\) be a short exact sequence
  (\cref{mathlib:short_exact}) of \(\cO_X\)-modules with \(\mathcal G\) and
  \(\mathcal H\) locally free (\cref{found:locally_free}); by
  \cref{found:isLocallyFree_kernel_of_shortExact} the kernel \(\mathcal K\) is then
  locally free as well.  Its rank satisfies
  \[
    \operatorname{rank}\mathcal K=\operatorname{rank}\mathcal G-\operatorname{rank}\mathcal H .
  \]
\end{lemma}

\begin{proof}
  \uses{found:locally_free, found:isLocallyFree_kernel_of_shortExact, mathlib:short_exact}
  By \cref{found:isLocallyFree_kernel_of_shortExact} the kernel \(\mathcal K\) is
  locally free, so all three terms carry a well-defined constant rank.  Locally the
  sequence splits --- a free quotient makes it split --- giving
  \(\mathcal G\cong\mathcal K\oplus\mathcal H\) on each trivialising open; the rank
  of a finite free module is additive over direct sums, so
  \(\operatorname{rank}\mathcal G=\operatorname{rank}\mathcal K+\operatorname{rank}
  \mathcal H\), which rearranges to the stated identity.
\end{proof}

% SOURCE (rank $r = h^0(L)-1$, Lazarsfeld, A Sampling of Vector Bundle Techniques in the Study of Linear Series, \S1.1 (p.~505) and \S1.3 (p.~509); read from references/lazarsfeld-vbt.pdf):
% SOURCE QUOTE (\S1.1, rank formula): "where $r = r(L) = h^0(L) - 1$."
% SOURCE QUOTE (\S1.3, rank of $M_L$): "Then there is a surjective evaluation map of vector bundles $e_L : H^0(L) \otimes_{\mathbf{C}} \mathcal{O}_X \longrightarrow L$, and we put $M_L = \ker e_L$. Thus $M_L$ is a vector bundle on $X$ of rank $r = r(L)$, and by construction one has the basic exact sequence $(1.3.1)$ $0 \longrightarrow M_L \longrightarrow H^0(L) \otimes_{\mathbf{C}} \mathcal{O}_X \longrightarrow L \longrightarrow 0$."
\begin{lemma}[Evaluation kernel of a globally generated locally free module is locally free]
  \label{found:eval_kernel_locally_free}
  \lean{AlgebraicGeometry.Scheme.Modules.evalKernel_isLocallyFree}
  \uses{found:eval_kernel_bundle, found:eval_shortexact_of_globally_generated, found:trivial_bundle_locally_free, found:isLocallyFree_kernel_of_shortExact, found:rank_kernel_of_shortExact, found:locally_free}
  \textit{Source: Kemeny, \S0 (preliminaries on locally free sheaves); rank formula after Lazarsfeld, VBT \S1.1--1.3.}
  If \(\mathcal F\) is globally generated --- so \(\mathrm{ev}_{\mathcal F}\) is
  surjective --- and locally free of finite rank \(r\), then the evaluation kernel
  \(M_{\mathcal F}\) (\cref{found:eval_kernel_bundle}) is locally free of rank
  \(\dim_k H^0(X,\mathcal F)-r\).  In particular, for a globally generated line bundle
  \(L\) (\(r=1\)), \(M_L\) is locally free of rank \(h^0(L)-1\).
\end{lemma}

\begin{proof}
  \uses{found:eval_kernel_bundle, found:eval_shortexact_of_globally_generated, found:trivial_bundle_locally_free, found:isLocallyFree_kernel_of_shortExact, found:rank_kernel_of_shortExact, found:locally_free}
  Global generation gives the short exact sequence
  \(0\to M_{\mathcal F}\to H^0(X,\mathcal F)\otimes_k\cO_X\to\mathcal F\to 0\)
  (\cref{found:eval_shortexact_of_globally_generated}), with middle term locally free
  of rank \(\dim_k H^0(X,\mathcal F)\) (\cref{found:trivial_bundle_locally_free}) and
  quotient \(\mathcal F\) locally free of rank \(r\) by hypothesis.  Hence the kernel
  \(M_{\mathcal F}\) is locally free (\cref{found:isLocallyFree_kernel_of_shortExact},
  \cref{found:locally_free}) of rank
  \(\dim_k H^0(X,\mathcal F)-r\) (\cref{found:rank_kernel_of_shortExact}); for a line
  bundle \(r=1\) this is \(h^0(L)-1\).
\end{proof}
